%%%%%%%%%%%%%%%%%%%%%%%%%%%%%%%%%%%%%%%%%%%%%%%%%%%%%%%%%%%%%%%%%%%%%%
% Lớp 12 | Tọa độ của véc-tơ trong không gian | Tọa độ của véc-tơ và các công thức
% Dạng 2H2?2-3: Tìm tọa độ véc-tơ
% Số câu: 79  (sinh tự động bởi tach_cau.py; từ nay thầy sửa trực tiếp file này)
%%%%%%%%%%%%%%%%%%%%%%%%%%%%%%%%%%%%%%%%%%%%%%%%%%%%%%%%%%%%%%%%%%%%%%

%%%%%%%%%%%%%%%% TRẮC NGHIỆM 4 PHƯƠNG ÁN %%%%%%%%%%%%%%%%

% nguồn: Lop_12_Hinh_hoc_Chuong_2
\begin{ex}%[2H2N2-3]%[ma:c4718362]
	Trong không gian $Oxyz$, cho $A(2;-1;0)$ và $B(1;1;-3)$. Véc-tơ $\overrightarrow{AB}$ có tọa độ là 
	\choice 
		 {$(-1;-2;3)$}
		{\True $(-1;2;-3)$} 
		{$(1;-2;3)$} 
		 {$(-1;-2;-3)$}
	 \loigiai{
	 $\overrightarrow{AB}=(1-2;1+1;-3-0)=(-1;2;-3)$.}
\end{ex}

% nguồn: Lop_12_Hinh_hoc_Chuong_2
\begin{ex}%[2H2N2-3]%[ma:6dd4bfba]
	Trong không gian $Oxyz$, cho hai điểm $A(1;-1;2)$ và $B(2;1;-4)$. Véctơ  $\overrightarrow{AB}$ có tọa độ
	\choice 
		 {$(3;0;-2)$}
		{$(-1;-2;6)$} 
		 {$(1;0;-6)$}
		{\True  $(1;2;-6)$}
	 \loigiai{
		 Ta có $\overrightarrow{AB}=(2-1;1-(-1);-4-2)=(1;2;-6)$.
	 }
\end{ex}

% nguồn: Lop_12_Hinh_hoc_Chuong_2
\begin{ex}%[2H2N2-3]%[ma:85f22243]
	Trong không gian với  hệ tọa độ $Oxyz$, cho vectơ $\overrightarrow{u} = -20 \overrightarrow{i} + 6 \overrightarrow{j} + 9 \overrightarrow{k}$. Tọa độ của vectơ $\overrightarrow{u}$ là
	 \choice
		{\True $(-20;6;9)$} 
		{$(20;6;9)$} 
		 {$(-20;6;0)$}
		 {$(20;6;0)$}
	 \loigiai{
		Do $\overrightarrow{u} = -20 \overrightarrow{i} + 6 \overrightarrow{j}  + 9 \overrightarrow{k}$ nên tọa độ vectơ $\overrightarrow{u} = (-20;6;9)$.
	 }
\end{ex}

% nguồn: Lop_12_Hinh_hoc_Chuong_2
\begin{ex}%[2H2N2-3]%[ma:cfe492d6]
	Trong không gian $Oxyz$, cho $\overrightarrow{a}=(2;3;2)$ và $\overrightarrow{b}=(1;1-1)$. Vectơ $\overrightarrow{a}-\overrightarrow{b}$ có tọa độ là 
	\choice 
		{$(3;4;1)$} 
		 {$(-1;-2;3)$}
		{$(3;5;1)$} 
		 {\True $(1;2;3)$}
	\loigiai{ 
		Vectơ $\overrightarrow{a}-\overrightarrow{b}$  có tọa độ $(1;2;3)$.
	} 
\end{ex}

% nguồn: Lop_12_Hinh_hoc_Chuong_2
\begin{ex}%[2H2N2-3]%[ma:fdf72ed7]
	Trong không gian $Oxyz$, cho  hai vectơ $\vec{u}=(3;2;-1)$, $\vec{v}=(5;-4;2)$. Toạ độ của vectơ $\vec{u}-\vec{v}$~là
	 \choice
		{\True  $(-2;6;-3)$}
		 {$(2;-6;3)$}
		 {$(-2;-2;-3)$}
		 {$(2;2;1)$}
	 \loigiai{
		Ta  có $\vec{u}-\vec{v}=(-2;6;-3)$.
	} 
\end{ex}

% nguồn: Lop_12_Hinh_hoc_Chuong_2
\begin{ex}%[2H2N2-3]%[ma:69633035]
	Trong không gian $Oxyz$, cho $\vec{u}=\vec{i}-3\vec{k}$.  Tọa độ của vectơ $\vec{u}$ là
	 \choice
		{$(1;-3;  0)$}
		{$(0;  2;-3)$}
		 {\True $(1; 0;-3)$}
		{$(1;-1; 3)$} 
	\loigiai{ 
		Do $\vec{u}=\vec{i}-3 \vec{k} \Rightarrow \vec{u}=1 \cdot \vec{i}+0 \cdot \vec{j}-3 \cdot  \vec{k} \Rightarrow \vec{u}=(1;0;-3)$.
	 }
\end{ex}

% nguồn: Lop_12_Hinh_hoc_Chuong_2
\begin{ex}%[2H2N2-3]%[ma:5d667281]
	Trong không gian với hệ  toạ độ $Oxyz$, tọa độ của vectơ $\overrightarrow{u}=3 \overrightarrow{i}-\overrightarrow{j}-2 \overrightarrow{k}$ là
	 \choice
		 {\True $(3;-1;-2)$}
		 {$(3;1;-2)$}
		 {$(3;-1;2)$}
		{$(3;1;2)$} 
	 \loigiai{
		Ta có  $\overrightarrow{u}=3 \overrightarrow{i}-\overrightarrow{j}-2 \overrightarrow{k} \Rightarrow \overrightarrow{u}=(3;-1;-2)$.
	 }
\end{ex}

% nguồn: Lop_12_Hinh_hoc_Chuong_2
\begin{ex}%[2H2N2-3]%[ma:e4d6fdfd]
	Trong không gian $Oxyz$, cho vectơ $\vec{u}=(1;-2;3)$.  toạ độ của vectơ $-3\vec{u}$~là
	 \choice
		 {$(3;-6;9)$}
		{$(-3;-6;-9)$} 
		 {$(3;6;9)$}
		{\True $(-3;6;-9)$} 
	 \loigiai{
		 Ta có $-3\vec{u}=(-3;6;-9)$.
	} 
\end{ex}

% nguồn: Lop_12_Hinh_hoc_Chuong_2
\begin{ex}%[2H2N2-3]%[ma:0619ab40]
	Trong không gian với hệ trục  tọa độ $Oxyz$, cho $\overrightarrow{a}=-\overrightarrow{i}+2\overrightarrow{j}-3\overrightarrow{k}$. Tọa độ của $\overrightarrow{a}$ là
	\choice 
		{$(-2;-1;-3)$} 
		 {$(-3;2;-1)$}
		{$(2;-3;-1)$} 
		{\True $(-1;2;-3)$} 
	\loigiai{ 
		Do đó,  $\overrightarrow{a}=-\overrightarrow{i}+2\overrightarrow{j}-3\overrightarrow{k}=(-1;2;-3)$.
	} 
\end{ex}

% nguồn: Lop_12_Hinh_hoc_Chuong_2
\begin{ex}%[2H2N2-3]%[ma:36cdc37b]
	Trong không gian $Oxyz$,  biết điểm $M(1;-2; 2)$. Tính $\overrightarrow{OM}$ theo các vectơ $\vec{i}$; $\vec{j}$; $\vec{k}$.
	\choice 
		{$\overrightarrow{OM}=-\vec{i} +2\vec{j}  -2\vec{k}$}
		 {$\overrightarrow{OM}= 2\vec{i} +3\vec{j} -4\vec{k}$}
		{$\overrightarrow{OM}=\vec{i} +2\vec{j} -2\vec{k}$} 
		{\True $\overrightarrow{OM}=\vec{i} -2\vec{j} +2\vec{k}$} 
	\loigiai{ 
		Ta có  $\overrightarrow{OM}=(1;-2; 2)$ nên $\overrightarrow{OM}=\vec{i} -2\vec{j} +2\vec{k}$.
	} 
\end{ex}

% nguồn: Lop_12_Hinh_hoc_Chuong_2
\begin{ex}%[2H2N2-3]%[ma:dae51991]
	Trong không  gian với hệ tọa độ $Oxyz$, cho vectơ $\vec{u}=-20\vec{i}+6\vec{j}+9\vec{k}$. Tọa độ của vectơ $\vec{u}$ là
	\choice 
		 {\True$(-20;6;9)$}
		 {$(20;6;9)$}
		{$(-20;6;0)$} 
		 {$(20;6;0)$}
	 \loigiai{
		Do $\vec{u}=-20\vec{i}+6\vec{j}+9\vec{k}$ nên tọa độ  vectơ $\vec{u}=(-20;6;9)$.
	} 
\end{ex}

% nguồn: Lop_12_Hinh_hoc_Chuong_2
\begin{ex}%[2H2N2-3]%[ma:2f745224]
	Trong không gian $Oxyz$, cho hai vectơ $\overrightarrow{a}=(2;-2;1)$, $\overrightarrow{b}=(0;1;3)$. Tọa độ của vectơ  $\overrightarrow{a}+\overrightarrow{b}$ là
	 \choice
		{$(2;-3;2)$} 
		 {$\left(1;-\dfrac{1}{2};2\right)$}
		{\True $(2;-1;4)$} 
		 {$(-2;3;2)$}
	 \loigiai{
		Ta có $\overrightarrow{a}+\overrightarrow{b}=(2;-1;4)$. 
	} 
\end{ex}

% nguồn: Lop_12_Hinh_hoc_Chuong_2
\begin{ex}%[2H2N2-3]%[ma:05d3cbe4]
	Trong không  gian $Oxyz$, cho hai vectơ $\overrightarrow{a} = (-2;5;2)$ và $\overrightarrow{v} = (1;-3;-1)$, vectơ $\overrightarrow{a} - \overrightarrow{b}$ có tọa độ là
	 \choice
		 {$(1;-2;3)$}
		 {\True $(-3;8;3)$}
		 {$(-1;2;1)$}
		{$(3;-8;-3)$} 
	\loigiai{ 
		Ta có $\overrightarrow{a} -  \overrightarrow{b} = (-2-1; 5-(-3); 2-(-1)) = (-3;8;3)$.
	} 
\end{ex}

% nguồn: Lop_12_Hinh_hoc_Chuong_2
\begin{ex}%[2H2N2-3]%[ma:4434cc16]
	Trong không gian $Oxyz$, tọa~độ của  vectơ $\vec{k}$~là
	\choice{$(1; 1; 1)$} 
		{$(1; 0;  0)$}
		{$(0;  1; 0)$}
		 {\True $(0; 0; 1)$}
	 \loigiai{
		 Ta có $\vec{k}=(0;0;1)$.
	 }
\end{ex}

% nguồn: Lop_12_Hinh_hoc_Chuong_2
\begin{ex}%[2H2H2-3]%[ma:bc13bbed]
	Trong không gian $Oxyz$, cho điểm $A(0;-2;-1)$  và $B(1;-1;2)$. Tìm tọa độ điểm $M$ sao cho $\overrightarrow{MA}=2\overrightarrow{MB}$.
	 \choice
		 {$M\left(\dfrac{2}{3};-\dfrac{4}{3};1\right)$}
		 {$M\left(\dfrac{1}{2};-\dfrac{3}{2};\dfrac{1}{2}\right)$}
		 {\True $M(2;0;5)$}
		 {$M(-1;-3;-4)$}
	\loigiai{ 
		Giả  sử $M(x;y;z)$.
		Ta có  $\overrightarrow{MA}=(0-x;-2-y;-1-z)$ và $\overrightarrow{MB}=(1-x;-1-y;2-z)$.\\
		Điều kiện $\overrightarrow{MA} = 2\overrightarrow{MB}$ tương đương  với hệ phương trình
		$$\heva{&-x=2\cdot(1-x)\\&-2-y=2\cdot(-1-y)\\&-1-z=2\cdot(2-z)} \Leftrightarrow \heva{&-x=2-2x\\&-2-y=-2-2y\\&-1-z=4-2z}  \Leftrightarrow \heva{&x=2\\&y=0\\&z=5.}$$
		 Vậy $M(2;0;5)$.
	 }
\end{ex}

% nguồn: Lop_12_Hinh_hoc_Chuong_2
\begin{ex}%[2H2H2-3]%[ma:7829ad62]
	Trong không gian $Oxyz$, cho  hai vectơ $\overrightarrow{a}=(-1;2x-1;1-3z)$ và $\overrightarrow{b}=(2+3y;-1;-2)$. Khi $\overrightarrow{a}=\overrightarrow{b}$thì tổng $T=x+2y^2+3z^3$ bằng
	\choice 
		{$2$} 
		{\True  $5$}
		 {$1$}
		{$4$} 
	\loigiai{ 
		 Ta có
		$$\overrightarrow{a}=\overrightarrow{b} \Leftrightarrow \heva{& -1=2+3y \\& 2x-1=-1 \\& 1-3z=-2} \Leftrightarrow \heva{& 3y=-3 \\& 2x=0\\& -3z=-3} \Leftrightarrow \heva{& y=-1\\&x=0\\&z=1.}$$ 
		Thay các  giá trị $x$, $y$, $z$ vào biểu thức $T$ ta được
		$$T=x+2y^2+3z^3=0+2(-1)^2+3(1)^3=0+2\cdot 1+3\cdot 1=5.$$ 
	 }
\end{ex}

% nguồn: Lop_12_Hinh_hoc_Chuong_2
\begin{ex}%[2H2H2-3]%[ma:ddc66d90]
	Trong không gian $Oxyz$, cho $=2\overrightarrow{i}+5\overrightarrow{j}-\overrightarrow{k},\,\overrightarrow{ON}=3\overrightarrow{i}-2\overrightarrow{j}$. Tìm tọa độ của  $\overrightarrow{MN}$.
	 \choice
		 {$\overrightarrow{MN}=(-1;-7;1)$}
		{\True $\overrightarrow{MN}=(1;-7;1)$} 
		 {$\overrightarrow{MN}=(1;-7;-1)$}
		 {$\overrightarrow{MN}=(1;7;1)$}
	\loigiai{ 
		Ta có $=2\overrightarrow{i}+5\overrightarrow{j}-\overrightarrow{k}\Rightarrow M(  2;5;-1)$ và $\overrightarrow{ON}=3\overrightarrow{i}-2\overrightarrow{j}\Rightarrow N( 3;-2;0)$.\\
		Do  đó $\overrightarrow{MN}=(1;-7;1)$.
	} 
\end{ex}

% nguồn: Lop_12_Hinh_hoc_Chuong_2
\begin{ex}%[2H2H2-3]%[ma:ef999721]
	Trong không gian $Oxyz$, cho $\vec{u}=\vec{i}-2 \vec{k}$.  Toạ độ của vectơ $\vec{u}$~là
	 \choice{$(1;-2; 0)$}
		{$(1; 0;  2)$}
		 {$(1;-2)$}
		{\True $(1;  0;-2)$}
	 \loigiai{
		 Ta có $\vec{u}=\vec{i}-2 \vec{k}=(1-2\cdot 0;0-2\cdot 0;0 -2\cdot 1)=(1;0;-2)$.
	} 
\end{ex}

% nguồn: Lop_12_Hinh_hoc_Chuong_2
\begin{ex}%[2H2H2-3]%[ma:a0b66d59]
	Trong không gian với hệ tọa độ $(O;\overrightarrow{i};\overrightarrow{j};\overrightarrow{k})$, cho hai vectơ  $\overrightarrow{a}=(1;2;3)$ và $\overrightarrow{b}=2\overrightarrow{i}-4\overrightarrow{k}$. Tính tọa độ vectơ $\overrightarrow{u}=\overrightarrow{a}-\overrightarrow{b}$
	 \choice
		 {$\overrightarrow{u}=(-1;2;-1)$}
		{$\overrightarrow{u}=(-1;-2;3)$} 
		 {$\overrightarrow{u}=(-1;6;3)$}
		 {\True $\overrightarrow{u}=(-1;2;7)$}
	\loigiai{ 
		Ta có $\overrightarrow{i}=(1;0;0)$,  $\overrightarrow{j}=(0;1;0)$, $\overrightarrow{k}=(0;0;1)$ nên $\overrightarrow{b}=2\overrightarrow{i}-4\overrightarrow{k}\Rightarrow \overrightarrow{b}=(2;0;-4)$.\\
		Suy ra  $\overrightarrow{u}=\overrightarrow{a}-\overrightarrow{b}=(-1;2;7)$.
	 }
\end{ex}

% nguồn: Lop_12_Hinh_hoc_Chuong_2
\begin{ex}%[2H2H2-3]%[ma:78b228d0]
	Trong không gian $Oxyz$,  toạ độ của vectơ $\vec{u}=\vec{j}-\vec{k}$~là
	\choice{$(1;-1; 0)$} 
		{$(1;  0;-1)$}
		{\True  $(0; 1;-1)$}
		{$(-1; 1;  0)$}
	\loigiai{ 
		Ta  có $\vec{j}-\vec{k}=(0-0;1-0;0-1)=(0;1;-1)$.
	 }
\end{ex}

% nguồn: Lop_12_Hinh_hoc_Chuong_2
\begin{ex}%[2H2H2-3]%[ma:e92ba7bd]
	Trong không gian $Oxyz$, cho vectơ $\overrightarrow{a}=(1;1;-3)$, $\overrightarrow{b}=(2;2;-2)$, $\overrightarrow{c}=2\overrightarrow{i}+2\overrightarrow{j}-6\overrightarrow{k}$ và $\overrightarrow{d}=(1;1;-1)$. Cặp  vectơ nào sau đây cùng phương?
	 \choice
		{$\overrightarrow{a}$ và  $\overrightarrow{b}$}
		 {$\overrightarrow{a}$ và $\overrightarrow{d}$}
		{\True $\overrightarrow{a}$ và  $\overrightarrow{c}$}
		{$\overrightarrow{b}$ và $\overrightarrow{c}$} 
	\loigiai{ 
		Tọa độ các vectơ $\overrightarrow{a}=(1;1;-3)$, $\overrightarrow{b}=(2;2;-2)$,  $\overrightarrow{c}=(2;2;-6)$, $\overrightarrow{d}=(1;1;-1)$.
		\begin{itemize}
			\item Xét cặp $\overrightarrow{a}$ và $\overrightarrow{b}$, ta  có $\dfrac{1}{2} = \dfrac{1}{2} \ne \dfrac{-3}{-2}$. Suy ra $\overrightarrow{a}$ và $\overrightarrow{b}$ không cùng phương.
			\item Xét cặp $\overrightarrow{a}$ và $\overrightarrow{d}$, ta có $\dfrac{1}{1} = \dfrac{1}{1} \ne \dfrac{-3}{-1}$. Suy ra $\overrightarrow{a}$ và  $\overrightarrow{b}$ không cùng phương.
			\item Xét cặp $\overrightarrow{a}$ và $\overrightarrow{c}$, ta có $\overrightarrow{c}=(2;2;-6)=2\cdot(1;1;-3)=2\cdot  \overrightarrow{a}$. Vậy $\overrightarrow{a}$ và $\overrightarrow{c}$ cùng phương.
			\item Xét cặp $\overrightarrow{b}$  và $\overrightarrow{c}$, ta có $\dfrac{2}{2} = \dfrac{2}{2} \ne \dfrac{-2}{-6}$. Suy ra $\overrightarrow{a}$ và $\overrightarrow{b}$ không cùng phương.
		\end{itemize}
	} 
\end{ex}

% nguồn: Lop_12_Hinh_hoc_Chuong_2
\begin{ex}%[2H2H2-3]%[ma:16c59368]
	 Trong không gian với hệ tọa độ $Oxyz$, cho hai véctơ là $\overrightarrow{u}=(2;1;-1)$ và $\overrightarrow{v}=(1;3;1)$. Tọa độ của véctơ $\overrightarrow{u}+2\overrightarrow{v}$ tương ứng là
	 \choice
		{$(3;4;0)$} 
		{$(1;-2;-2)$} 
		 {\True $(4;7;1)$}
		{$(5;5;-1)$} 
	 \loigiai{
		 Ta có $\overrightarrow{u}+2\overrightarrow{v}=(4;7;1)$.
	} 
\end{ex}

% nguồn: Lop_12_Hinh_hoc_Chuong_2
\begin{ex}%[2H2H2-3]%[ma:e6412af6]
	Trong không gian $Oxyz$,  cho ba vectơ $\overrightarrow{a}=(3;4;2)$, $\overrightarrow{b}=(-5;0;3)$ và $\overrightarrow{c}=(1;2;-4)$. Tìm tọa độ của vectơ $\overrightarrow{u} =3\overrightarrow{a}+2\overrightarrow{b}-\overrightarrow{c}$.
	\choice 
		{\True  $\overrightarrow{u}=(-2;10;16)$}
		 {$\overrightarrow{u}=(2;10;-16)$}
		 {$\overrightarrow{u}=(-1;5;8)$}
		{$\overrightarrow{u}=(-2;-10;16)$} 
	 \loigiai{
		Ta  có
		$3\overrightarrow{a}=3\cdot (3;4;2)=(9;12;6)$, $2\overrightarrow{b}=2\cdot(-5;0;3)=(-10;0;6)$,  $-\overrightarrow{c}=-(1;2;-4)=(-1;-2;4)$.\\
		Khi đó  $\overrightarrow{u}=3\overrightarrow{a}+2\overrightarrow{b}-\overrightarrow{c}=(9-10-1;12+0-2;6+6+4)=(-2;10;16)$.
	} 
\end{ex}

% nguồn: Lop_12_Hinh_hoc_Chuong_2
\begin{ex}%[2H2H2-3]%[ma:b08fc057]
	Trong không gian $Oxyz$, cho  biểu diễn của vectơ $\overrightarrow{a}$ qua các vectơ đơn vị là $\overrightarrow{a}=2\overrightarrow{i}+\overrightarrow{k}-3\overrightarrow{j}$. Tọa độ của vectơ $\overrightarrow{a}$ là
	\choice 
		{\True $(2;-3;1)$} 
		 {(1;-3;2)}
		 {$(2;1;-3)$}
		{$(1;2;-3)$} 
	\loigiai{ 
		Ta  có $\overrightarrow{a}=2\overrightarrow{i}+\overrightarrow{k}-3\overrightarrow{j}=2\overrightarrow{i}-3\overrightarrow{j}+\overrightarrow{k}\Rightarrow \overrightarrow{a}=(2;-3;1)$.
	} 
\end{ex}

% nguồn: Lop_12_Hinh_hoc_Chuong_2
\begin{ex}%[2H2H2-3]%[ma:39dfc461]
	Trong mặt  phẳng $Oxyz$, cho $\overrightarrow{a}=(3;-1;2)$, $\overrightarrow{b}=(4;2;-6)$. Giá trị của $\left|\overrightarrow{a}+\overrightarrow{b}\right|$ bằng
	 \choice
		 {$66$}
		{\True  $\sqrt{66}$}
		{$3 \sqrt{14}$} 
		{$2$} 
	 \loigiai{
		Ta có  $$\overrightarrow{a}+\overrightarrow{b}=(7;1;-4)\Rightarrow\left|\overrightarrow{a}+\overrightarrow{b}\right|=\sqrt{7^2+1^2+(-4)^2}=\sqrt{66}.$$
	} 
\end{ex}

% nguồn: Lop_12_Hinh_hoc_Chuong_2
\begin{ex}%[2H2H2-3]%[ma:d1b1c0d2]
	Cho hình chóp tứ giác đều $S.ABCD$ có cạnh đáy bằng $a\sqrt{2}$, cạnh bên bằng $a\sqrt{5}$. Gọi $O$ là tâm của hình vuông $ABCD$. Với hệ toạ độ $Oxyz$ được thiết lập như hình bên (gốc tọa độ $O$ trùng  với tâm hình vuông $ABCD$), tọa độ $\overrightarrow{SC}$ là
	\begin{center}
		\begin{tikzpicture}[scale=0.8,font=\footnotesize,line join=round,line cap=round,>=stealth]
				\coordinate (A) at (0,0); 
				\coordinate  (B) at (-2,-2);
				 \coordinate (D) at (5,0);
				\coordinate (C)  at ($(B)+(D)-(A)$);
				 \coordinate (O) at ($(A)!0.5!(C)$);
				\coordinate (S) at ($(O)+(0,5)$); 
				\foreach \i  in {B,C,D}{\draw (S)--(\i);}
				 \draw (B)--(C)--(D);
				 \draw[dashed,thin](S)--(A)--(B)--(D)--(A)--(C);
				\draw (S)--(B) 
					(S)--(C) 
					(S)--(D) 
					(B)--(C)--(D); 
				 \draw[dashed,thin](S)--(A)--(C)
					 (A)--(B)
					 (A)--(D)--(B)
					(S)--(O); 
				\pic[draw,thin,angle radius=2mm] {right angle =  C--O--D};
				\foreach \i/\g in {S/60,A/-90,B/-90,C/-90,D/-90,O/-90}{\draw (\i)  circle (1.5pt)
					 ($(\i)+(\g:4mm)$) node[scale=1]{$\i$};}
				\coordinate (Z) at ($(O)!1.2!(S)$); 
				\coordinate (Y) at ($(O)!1.2!(D)$); 
				\coordinate (X) at ($(O)!2!(C)$); 
				\draw[->]  (S)--(Z);
				 \draw[->] (C)--(X);
				\draw[->]  (D)--(Y);
				\draw (X) node[scale=1,left]{$x$}; 
				\draw (Y) node[scale=1,above]{$y$}; 
				\draw  (Z) node[scale=1,left]{$z$};
			\end{tikzpicture}
	\end{center}
	\choice 
		{$\overrightarrow{SC}=\left(2a; 0;-2a\right)$} 
		 {$\overrightarrow{SC}=\left(2a;-a;-2a\right)$}
		{\True $\overrightarrow{SC}=\left(a; 0;-2a\right)$} 
		{$\overrightarrow{SC}=\left(a;  0; 2a\right)$}
	\loigiai{ 
		Ta có tọa độ điểm  $C=(a;0;0)$ và $S=(0;0;2a)$;\\
		Tọa độ $\overrightarrow{SC}$ là $\overrightarrow{SC}=\left(a;  0;-2a\right)$.
	 }
\end{ex}

% nguồn: Lop_12_Hinh_hoc_Chuong_2
\begin{ex}%[2H2H2-3]%[ma:62c27e2f]
	Trong không gian $Oxyz$, cho ba điểm $A(2;-1;5)$, $B(5;-5;7)$; $M(x;y;1)$. Khi $A$, $B$, $M$ thẳng hàng  thì giá trị của $x;y$ là
	 \choice
		{$x=4$;  $y=-7$}
		 {\True $x=-4$; $y=7$}
		{$x=4$; $y=7$} 
		{$x=-4$;  $y=-7$}
	 \loigiai{
		Ta có  $\overrightarrow{AB}=(3;-4;2)$, $\overrightarrow{AM}=(x-2;y+1;-4)$\\
		Để ba điểm $A$,  $B$, $M$ thẳng hàng thì $\dfrac{x-2}{3}=\dfrac{y+1}{-4}=\dfrac{-4}{2} \Rightarrow \heva{&x=-4 \\& y=7.}$
	 }
\end{ex}

% nguồn: Lop_12_Hinh_hoc_Chuong_2
\begin{ex}%[2H2H2-3]%[ma:3b4ccf01]
	Trong không gian $Oxyz$, toạ độ  của vectơ $\vec{u}=\vec{i}-\vec{j}$~là
	 \choice{\True $(1;-1; 0)$}
		{$(1;  0;-1)$}
		 {$(0; 1;-1)$}
		 {$(-1; 1; 0)$}
	 \loigiai{
		Ta có $\vec{i}-\vec{j}=(1-0;0-1;0-0)=(1;-1;0)$. 
	 }
\end{ex}

% nguồn: Lop_12_Hinh_hoc_Chuong_2
\begin{ex}%[2H2H2-3]%[ma:b537e708]
	Trong không gian $Oxyz$, cho $\overrightarrow{a}=(-3;5;2)$,  $\overrightarrow{b}=(0;-1;3)$, $\overrightarrow{c}=(1;-1;1)$. Tọa độ $\overrightarrow{v}=2\overrightarrow{a}-3\overrightarrow{b}+15\overrightarrow{c}$ là
	 \choice
		{$\overrightarrow{v}=(-9;2;10)$} 
		{$\overrightarrow{v}=(9;-1;10)$} 
		{$\overrightarrow{v}=(9;2;10)$} 
		{\True $\overrightarrow{v}=(9;-2;10)$} 
	 \loigiai{
		Ta có  $2\overrightarrow{a}=(-6;10;4)$, $3\overrightarrow{b}=(0;-3;9)$, $15\overrightarrow{c}=(15;-15;15)$.\\
		Suy  ra $\overrightarrow{v}=2\overrightarrow{a}-3\overrightarrow{b}+15\overrightarrow{c}=(9;-2;10)$.
	 }
\end{ex}

% nguồn: Lop_12_Hinh_hoc_Chuong_2
\begin{ex}%[2H2H2-3]%[ma:5f55b01d]
	 Trong không gian với hệ tọa độ $Oxyz$, cho ba điểm $A(-1;1;2)$, $B(0;1;-1)$ và $C(x+2;y;-2)$ thẳng hàng. Tổng $x+y$ bằng
	 \choice
		{$\dfrac{7}{3}$} 
		 {$-\dfrac{8}{3}$}
		 {\True$-\dfrac{2}{3}$}
		{$-\dfrac{1}{3}$} 
	 \loigiai{
		Ta có $\overrightarrow{AB}=(1;0;-3)$,  $\overrightarrow{AC}=(x+3;y-1;-4)$.\\
		Vì $A$, $B$, $C$ thẳng hàng nên $\overrightarrow{AB}$ cùng  phương với $\overrightarrow{AC}$
		$$\Leftrightarrow\heva{&\dfrac{x+3}{1}=\dfrac{-4}{-3}\\&y-1=0}\Leftrightarrow\heva{&x=-\dfrac{5}{3}\\&y=1.}$$ 
		Vậy  $x+y=\dfrac{5}{3}+1=-\dfrac{2}{3}$.
	 }
\end{ex}

% nguồn: Lop_12_Hinh_hoc_Chuong_2
\begin{ex}%[2H2H2-3]%[ma:888497ab]
	Trong không gian $Oxyz$, cho $\vec{a}=(2;-3;  3)$, $\vec{b}=(0; 2;-1)$, $\vec{c}=(3;-1; 5)$. Tìm tọa độ của vectơ $\vec{u}=2 \vec{a}+3 \vec{b}-2 \vec{c}$.
	\choice 
		 {$(10;-2; 13)$}
		 {\True $(-2; 2;-7)$}
		{$(-2;-2; 7)$} 
		 {$(-2; 2; 7)$}
	\loigiai{ 
		Ta có $2 \vec{a}=(4;-6; 6)$,  $3 \vec{b}=(0; 6;-3)$, $-2 \vec{c}=(-6; 2;-10)$.\\
		Do đó $\vec{u}=2 \vec{a}+3 \vec{b}-2 \vec{c}=(-2;  2;-7)$.
	} 
\end{ex}

% nguồn: Lop_12_Hinh_hoc_Chuong_2
\begin{ex}%[2H2H2-3]%[ma:78fa83dd]
	Trong không  gian $Oxyz$, cho hai vectơ $\vec{a}=(1; 0;-3)$, $\vec{b}=(-1;-2; 0)$. Vectơ $\left[\vec{b}, \vec{a}\right]$ có tọa~độ
	\choice{$(-6; 3;-2)$} 
		{\True  $(6;-3; 2)$}
		{$(-6; 3;  2)$}
		{$(6; 3;  2)$}
	\loigiai{ 
		Ta  có $\left[\vec{b}, \vec{a}\right]= (6;-3;2)$.
	} 
\end{ex}

% nguồn: Lop_12_Hinh_hoc_Chuong_2
\begin{ex}%[2H2H2-3]%[ma:770a4eec]
	Trong không gian với hệ trục tọa độ $Oxyz$ cho hai vector $\overrightarrow{a}(1;1;-2)$, $\overrightarrow{b}=(-2;1;4)$. Tìm tọa  độ của vector $\overrightarrow{u}=\overrightarrow{a}-2\overrightarrow{b}$.
	 \choice
		{\True  $(5;-1;-10)$}
		{$(0;3;0)$} 
		 {$(-3;3;6)$}
		 {$(5;-1;10)$}
	 \loigiai{
		Ta có $\heva{&\overrightarrow{a}=(1;1;-2)\\&\overrightarrow{b}=(-2;1;4)}\Rightarrow \overrightarrow{u}=\overrightarrow{a}-2\overrightarrow{b}=(5;-1;-10)$. 
	} 
\end{ex}

% nguồn: Lop_12_Hinh_hoc_Chuong_2
\begin{ex}%[2H2V2-3]%[ma:8391499e]
	Trong không gian $Oxyz$, cho tam giác $ABC$ có $A(1;-1;-2)$ và  trọng tâm $G(2;1;-3)$. Tọa độ của vectơ $\overrightarrow{u}=\overrightarrow{AB}+\overrightarrow{AC}$ là
	 \choice
		 {$(3;6;3)$}
		 {\True $(3;6;-3)$}
		 {$(3;-3;6)$}
		{$(3;2;1)$} 
	\loigiai{ 
		Gọi $M$ là trung điểm cạnh $BC$. Lúc đó 
		$$\overrightarrow{u}=\overrightarrow{AB}+\overrightarrow{AC}=2\cdot\overrightarrow{AM}=2\cdot\dfrac{3}{2}\cdot\overrightarrow{AG}=3\overrightarrow{AG}=(3;6;-3).$$ 
	 }
\end{ex}

% nguồn: Lop_12_Hinh_hoc_Chuong_2
\begin{ex}%[2H2B2-3]%[ma:3406ddfa]
	Trong không gian $Oxyz$, cho hai vectơ $\vec{a}=\left(6;0;-3 \right),\vec{b}=\left(-1;-2;0 \right)$. Vectơ  $\vec{d}=\vec{a}-\vec{b}$ có tọa~độ
	 \choice
		{$\left(5;-2;-3  \right)$}
		{$\left(-5;2;3  \right)$}
		 {\True $\left(7;2;-3 \right)$}
		 {$\left(-7;-2;3 \right)$}
	 \loigiai{
		 Ta có $\vec{c}=\vec{a}-\vec{b}=\left(6+1;0+2;-3-0\right)=\left(7;2;-3\right)$.
	 }
\end{ex}

% nguồn: Lop_12_Hinh_hoc_Chuong_2
\begin{ex}%[2H2B2-3]%[ma:1dbe29e3]
	Trong không gian $Oxyz$, cho hai vectơ $\vec{a}=\left(6;0;-3 \right),\vec{b}=\left(-1;-2;0 \right)$. Vectơ $\vec{e}=\vec{b}-\vec{a}$ có  tọa~độ
	\choice 
		 {$\left(5;-2;-3 \right)$}
		 {$\left(-5;2;3 \right)$}
		{$\left(7;2;-3  \right)$}
		 {\True $\left(-7;-2;3 \right)$}
	\loigiai{ 
		 Ta có $\vec{d}=\vec{b}-\vec{a}=\left(-1-6;-2-0;0+3\right)=\left(-7;-2;3\right)$.
	} 
\end{ex}

% nguồn: Lop_12_Hinh_hoc_Chuong_2
\begin{ex}%[2H2B2-3]%[ma:7e8d9509]
	Trong không  gian $Oxyz$, cho hai vectơ $\vec{a}=\left(6;0;-3 \right),\vec{b}=\left(-1;-2;0 \right)$. Vectơ $\vec{c}=\vec{a}+\vec{b}$ có tọa~độ
	 \choice
		 {\True $\left(5;-2;-3 \right)$}
		 {$\left(-5;2;3 \right)$}
		{$\left(7;2;-3 \right)$} 
		{$\left(-7;-2;3  \right)$}
	\loigiai{ 
		Ta có $\vec{c}=\vec{a}+\vec{b}=\left(6-1;0-2;-3+0\right)=\left(5;-2;-3\right)$. 
	} 
\end{ex}

% nguồn: Lop_12_Hinh_hoc_Chuong_2
\begin{ex}%[2H2B2-3]%[ma:f98170a2]
	Trong không gian $Oxyz$, cho hai  vectơ $\vec{a}=\left(6;0;-3 \right),\vec{b}=\left(-1;-2;0 \right)$. Vectơ $\vec{f}=-\vec{a}-\vec{b}$ có tọa~độ
	\choice 
		 {$\left(5;-2;-3 \right)$}
		{\True  $\left(-5;2;3 \right)$}
		{$\left(7;2;-3 \right)$} 
		{$\left(-7;-2;3 \right)$} 
	\loigiai{ 
		 Ta có $\vec{f}=-\vec{a}-\vec{b}=\left(-6+1;0+2;3-0\right)=\left(-5;2;3\right)$.
	 }
\end{ex}

% nguồn: Lop_12_Hinh_hoc_Chuong_2
\begin{ex}%[2H2B2-3]%[ma:136de626]
	Trong  không gian $Oxyz$, cho hai vectơ $\vec{a}=\left(1;2;-3 \right)$. Vectơ $-2\vec{a}$ có tọa~độ
	 \choice
		{$\left(2;4;-6  \right)$}
		{\True $\left(-2;-4;6  \right)$}
		 {$\left(2;-4;-6 \right)$}
		 {$\left(-2;4;-6 \right)$}
	\loigiai{ 
		Ta có $-2\vec{a}=\left(-2\cdot 1;-2\cdot 2;-2\cdot (-3)\right)=(-2;-4;6)$. 
	 }
\end{ex}

% nguồn: Lop_12_Hinh_hoc_Chuong_2
\begin{ex}%[2H2B2-3]%[ma:4d3b7a75]
	Trong không gian  $Oxyz$, cho hai vectơ $\vec{a}=\left(1;2;-3 \right)$. Vectơ $2\vec{a}$ có tọa~độ
	\choice 
		{\True  $\left(2;4;-6 \right)$}
		{$\left(-2;-4;6  \right)$}
		 {$\left(2;-4;-6 \right)$}
		{$\left(-2;4;-6  \right)$}
	 \loigiai{
		Ta có  $2\vec{a}=\left(2\cdot 1;2\cdot 2;2\cdot (-3)\right)=(2;4;-6)$.
	 }
\end{ex}

% nguồn: Lop_12_Hinh_hoc_Chuong_2
\begin{ex}%[2H2B2-3]%[ma:3f7580b2]
	Trong không gian $Oxyz$, cho hai điểm $M\left(2;3;1  \right),N\left(3;1;5 \right)$. Tọa độ của vectơ $\overrightarrow{MN}$~là
	\choice 
		{$\left(-5;-4;-6 \right)$} 
		{$\left(5;4;6 \right)$} 
		{$\left(-1;2;-4 \right)$} 
		 {\True $\left(1;-2;4 \right)$}
	 \loigiai{
		Ta có $\overrightarrow{MN}=\left(3-2;1-3;5-1\right)=\left(1;-2;4\right)$. 
	} 
\end{ex}


%%%%%%%%%%%%%%%% ĐÚNG/SAI %%%%%%%%%%%%%%%%

% nguồn: Lop_12_Hinh_hoc_Chuong_2
\begin{ex}%[2H2N2-3]%[ma:5cf18e8c]
	Trong không gian  $Oxyz$, cho tam giác $ABC$ với $A(1; -1; 1)$, $B(2; 1; -2)$, $C(0; 0; 1)$.
	\choiceTF 
		{\True Toạ độ trọng tâm $G$ của tam giác $ABC$  là $(1; 0; 0)$}
		{\True Toạ độ điểm $E \in Ox$ sao cho $AE  \perp BC$ là $(-1; 0; 0)$}
		{Toạ độ điểm $M$ thoả mãn $\overrightarrow{AM} = \overrightarrow{AB}  - \overrightarrow{AC}$ là $(3; 0; 2)$}
		{\True Trực tâm $H$ của tam giác $ABC$ có tọa độ là $\left(  \dfrac{5}{9}; -\dfrac{4}{9}; \dfrac{8}{9} \right)$}
	 \loigiai{
		\begin{itemchoice}
			\itemch Trọng tâm $G$ có tọa độ\\ 
			$$G = \left( \dfrac{1 + 2 + 0}{3}; \dfrac{-1 + 1 + 0}{3}; \dfrac{1 -  2 + 1}{3} \right) = (1; 0; 0).$$
			\itemch Vì $E \in Ox \Rightarrow E(x; 0;  0)$. Ta có\\
			 $\overrightarrow{AE} = (x - 1; 1; -1)$, $\overrightarrow{BC} = (-2; -1; 3)$\\
			$\overrightarrow{AE} \cdot \overrightarrow{BC} = 0  \Leftrightarrow (x - 1)(-2) + 1 \cdot (-1) + (-1) \cdot 3 = 0 \Rightarrow -2x + 2 -1 -3 = 0 \Rightarrow x = -1$\\
			$\Rightarrow E  = (-1; 0; 0)$.
			\itemch Gọi $M(x; y; z) \Rightarrow \overrightarrow{AM} = (x - 1; y + 1; z -  1)$.\\
			Ta có $\overrightarrow{AB} = (1; 2; -3)$, $\overrightarrow{AC} =  (-1; 1; 0) \Rightarrow \overrightarrow{AB} - \overrightarrow{AC} = (2; 1; -3)$.\\
			$\Rightarrow \overrightarrow{AM} = \overrightarrow{AB} - \overrightarrow{AC} \Leftrightarrow \heva{ 
			&x  - 1 = 2 \\
			&y + 1 = 1  \\
			 &z - 1 = -3
			} \Rightarrow  \heva{
			&x = 3  \\
			&y  = 0 \\
			&z =  -2.
			} \Rightarrow M  = (3; 0; -2)$.
			\itemch Gọi $H(x; y; z)$ là trực  tâm. Ta có\\
			$\overrightarrow{BC} = (-2; -1; 3)$,  $\overrightarrow{AC} = (-1; 1; 0)$, $[\overrightarrow{BC}, \overrightarrow{AC}] = (-3; -3; -3)$\\
			$\overrightarrow{AH} =  (x - 1; y + 1; z - 1)$, $\overrightarrow{BH} = (x - 2; y - 1; z + 2)$, $\overrightarrow{CH} = (x; y; z - 1)$\\
			 Khi đó\\
			$\heva{ 
			&\overrightarrow{AH} \cdot \overrightarrow{BC} = 0 \\ 
			&\overrightarrow{BH} \cdot \overrightarrow{AC} =  0 \\
			&[\overrightarrow{BC}, \overrightarrow{AC}] \cdot \overrightarrow{CH}  = 0
			 }$
			$\Leftrightarrow  \heva{
			&- 2 ( x - 1) - ( y + 1)  + 3 ( z - 1)= 0\\
			&- ( x - 2) + ( y - 1)  = 0\\
			&- 3 x - 3 y - 3 ( z - 1)  = 0}
			 \Leftrightarrow \heva{
			&x=\dfrac{5}{9}  \\
			&y=-\dfrac{4}{9}  \\
			&z=\dfrac{8}{9}}.$\\ 
			Vậy $H = \left( \dfrac{5}{9}; -\dfrac{4}{9}; \dfrac{8}{9} \right)$. 
		\end{itemchoice}
	} 
\end{ex}

% nguồn: Lop_12_Hinh_hoc_Chuong_2
\begin{ex}%[2H2H2-3]%[ma:34621ec0]
	\immini[thm]{ 
	Cho  hình hộp $ABCD.A'B'C'D'$, biết điểm $A(0;0;0)$, $B(1;0;0)$, $C(1;2;0)$, $D'(-1;3;5)$. Gọi $M$, $N$ là tâm của các hình bình hành $ABB'A'$, $ADD'A'$.
	 \choiceTF
		{\True Tọa  độ $D(0;2;0)$}
		{\True Tọa  độ $A'(-1;1;5)$}
		{Tọa độ $\overrightarrow{MN}=(-1;  1;0)$}
		{$\left|\overrightarrow{AB}+\overrightarrow{AD}+\overrightarrow{CC'}\right|=\sqrt{29}$} 
		} 
		{ 
		\begin{tikzpicture}[scale=1.1,font=\footnotesize,line join=round,
					line  cap=round,>=stealth]
				 \def\a{2.3}
				\def\b{1} 
				\def\g{30} 
				\def\h{2} 
				\path (0:0) coordinate  (A)--++(\g:\b) coordinate (B)--++(0:\a) coordinate (C)--++(\g-180:\b) coordinate (D)
				\foreach  \x in {A,B,C,D}{($(\x)+(80:\h)$) coordinate (\x')}
					 ($(A)!.5!(B')$) coordinate (M)
					($(A)!.5!(D')$) coordinate  (N);
				\draw[dashed](A)--(B') 
					(A')--(B) 
					(B')--(B)--(A) 
					(B)--(C); 
				\draw (D)--(D') 
					(A)--(D') 
					(A')--(D) 
					 (A')--(B')--(C')--(D')--(A')--(A)--(D)--(C)--(C');
				 \foreach \x/\g in {A/180,B/-90,C/0,D/-30,A'/180,B'/180,C'/0,D'/-30,M/190,N/90}
				\fill (\x) circle (1pt) 
					($(\x)+(\g:2.3mm)$) node{$\x$}; 
			\end{tikzpicture}
		} 
	 \loigiai{
		\begin{itemchoice}
			\itemch Ta có  $\overrightarrow{AC}=(1;2;0)$, $\overrightarrow{AB}=(1;0;0)$.\\
			Theo qui tắc hình bình hành, ta  có $\overrightarrow{AD}=\overrightarrow{AC}-\overrightarrow{AB}=(0;2;0)$. \\
			Mà $A(0;0;0)$,  suy ra $D(0;2;0)$.
			\itemch Ta có $\overrightarrow{AA'}=\overrightarrow{DD'}=(-1;1;5)$, mà $A(0;0;0)$ nên có  $A'(-1;1;5)$.
			\itemch Ta có  $\overrightarrow{MN}=\dfrac{1}{2}\overrightarrow{BD}=\left(-\dfrac{1}{2};1;0\right)$.
			\begin{center}
				\begin{tikzpicture}[scale=1.1,font=\footnotesize,line join=round,
						line cap=round,>=stealth] 
						\def\a{2.3} 
						\def\b{1} 
						\def\g{30} 
						\def\h{2} 
						\path (0:0) coordinate (A)--++(\g:\b) coordinate  (B)--++(0:\a) coordinate (C)--++(\g-180:\b) coordinate (D)
						 \foreach \x in {A,B,C,D}{($(\x)+(80:\h)$) coordinate (\x')}
						($(A)!.5!(B')$)  coordinate (M)
						($(A)!.5!(D')$) coordinate (N); 
						\draw[dashed]  (M)--(N)
						(A)--(B') 
						(A')--(B) 
						 (B')--(B)--(A)
						(B)--(C) 
						 (B)--(D);
						\draw  (D)--(D')
						 (A)--(D')
						 (A')--(D)
						 (A')--(B')--(C')--(D')--(A')--(A)--(D)--(C)--(C');
						\foreach \x/\g in  {A/180,B/-90,C/0,D/-30,A'/180,B'/180,C'/0,D'/-30,M/190,N/90}
						\fill  (\x) circle (1pt)
						($(\x)+(\g:2.3mm)$) node{$\x$}; 
					\end{tikzpicture}
			\end{center}
			\itemch  Ta có $\overrightarrow{AC'}=\overrightarrow{AB}+\overrightarrow{AD}+\overrightarrow{AA'}=(0;3;5)$.\\
			Lúc đó $\left|\overrightarrow{AB}+\overrightarrow{AD}+\overrightarrow{CC'}\right|=\left|\overrightarrow{AB}+\overrightarrow{AD}+\overrightarrow{AA'}\right|=\left|\overrightarrow{AC'}\right|=\sqrt{0^2+3^2+5^2}=\sqrt{34}$. 
		\end{itemchoice}
	} 
\end{ex}

% nguồn: Lop_12_Hinh_hoc_Chuong_2
\begin{ex}%[2H2H2-3]%[ma:b039fe2c]
	 \immini{
	Trong không gian $Oxyz$, cho  hình hộp $OABC.O'A'B'C'$ có $A(1;1;-1)$, $B(0;3;0)$, $\overrightarrow{BC'}=(2;-6;6)$. Gọi $H$, $K$ lần lượt là trọng tâm của tam giác $OA'O'$ và $CB'C'$.
	 }
	{\begin{tikzpicture}[scale=0.5,font=\footnotesize,line join=round,line cap=round,>=stealth]
		\coordinate (O) at (0,0); 
		 \coordinate (A) at (-2,-2);
		\coordinate (C)  at (5,0);
		\coordinate  (B) at ($(A)+(C)-(O)$);
		\coordinate (H)  at (1,4);
		\foreach \i  in {O,A,B,C}{\coordinate (\i') at ($(\i)+(H)$);}
		\draw (O')--(A')--(B')--(C')--cycle; 
		 \draw (A)--(A')
			 (B)--(B')
			(C)--(C') 
			 (A)--(B)--(C);
		 \draw[dashed,thin](A)--(O)--(O')
			 (O)--(C);
		\foreach  \i/\g in {O'/90,A'/90,B'/90,C'/90,O/-90,A/-90,B/-90,C/-90}{\draw[fill=blue](\i) circle (1pt)
			($(\i)+(\g:5mm)$)  node[scale=1]{$\i$};}
	\end{tikzpicture}
	} 
	\choiceTF 
		{\True Tọa độ điểm  $C'$ là $(2;-3;6)$}
		{\True Tọa độ điểm $O'$  là $(3;-5;5)$}
		{Tọa độ véc to  $\overrightarrow{AB'}=(-2;3;-6)$}
		{Tọa độ véc tơ  $\overrightarrow{HK}=(-1;2;-1)$}
	 \loigiai{
		\begin{itemchoice}
			\itemch Gọi $C'(x;y;z)$. Ta  có $\overrightarrow{BC'}=(2;-6;6)\Rightarrow\heva{&x-0=2\\&y-3=-6\\&z-0=6}\Leftrightarrow\heva{&x=2\\&y=-3\\&z=6}\Rightarrow C'(2;-3;6)$.
			\itemch Gọi  $O'(x;y;z)$. Theo hình vẽ thì $\overrightarrow{AO'}=\overrightarrow{BC'}\Leftrightarrow\heva{&x-1=2\\&y-1=-6\\&z+1=6}\Leftrightarrow\heva{&x=3\\&y=-5\\&z=5}\Rightarrow O'(3;-5;5)$.
			\itemch Theo hình  vẽ thì $\overrightarrow{AB'}=\overrightarrow{OC'}=(2;-3;6)$.
			\itemch Ta có $\overrightarrow{HK}=\overrightarrow{AB}=(-1;2;1)$. 
		\end{itemchoice}
	} 
\end{ex}

% nguồn: Lop_12_Hinh_hoc_Chuong_2
\begin{ex}%[2H2H2-3]%[ma:04e924d7]
	Cho các điểm $A(1;-2;3)$, $B(-2;1;2)$,  $C(3;-1;2)$. Khi đó
	 \choiceTF
		{\True $\overrightarrow{AB}=(-3;3;-1)$} 
		{$\overrightarrow{AC}=(-2;-1;1)$} 
		{$\overrightarrow{AB}=3 \overrightarrow{AC}$} 
		{\True Ba  điểm $A$, $B$, $C$ không thẳng hàng}
	 \loigiai{
		\begin{itemchoice}
			\itemch Ta  có $\overrightarrow{AB}=\left(x_B-x_A;y_B-y_A;z_B-z_A\right)=(-3;3;-1)$.
			\itemch Ta có  $\overrightarrow{AC}=\left(x_C-x_A;y_C-y_A;z_C-z_A\right)=(2;1;-1)$.
			\itemch Ta có $\overrightarrow{AB}=(-3;3;-1)$, $\overrightarrow{AC}=(2;1;-1)$, $3\overrightarrow{AC}=(6;3;-3)\ne \overrightarrow{AB}$. 
			\itemch Vì $\dfrac{-3}{2}\ne \dfrac{3}{1}\ne \dfrac{-1}{-1}$ nên hai vectơ này không cùng  phương.\\
			Vậy  ba điểm $A$, $B$, $C$ không thẳng hàng.
		\end{itemchoice}
	} 
\end{ex}

% nguồn: Lop_12_Hinh_hoc_Chuong_2
\begin{ex}%[2H2H2-3]%[ma:f397dc60]
	Trong không gian $Oxyz$, cho $\overrightarrow{a}=\overrightarrow{i}+3\overrightarrow{k}-4\overrightarrow{j}$ và $\overrightarrow{b}=\left(m-n;4m-6n;n^{2}-3m+2\right)$, với $m,n$  là tham số.
	\choiceTF 
		{Tọa độ  $\overrightarrow{a}=(1;3;-4)$}
		{\True Dựng điểm  $A$ thỏa $\overrightarrow{OA}=\overrightarrow{a}$ thì $A(1;-4;3)$}
		{Tồn tại giá trị của $m$  và $n$ để $\overrightarrow{b}=\overrightarrow{0}$}
		{\True Nếu  $\overrightarrow{a}=\overrightarrow{b}$ thì $m+n=9$}
	\loigiai{ 
		\begin{itemchoice}
			\itemch  Tọa độ $\overrightarrow{a}=(1;-4;3)$.
			\itemch Khi $\overrightarrow{OA}=\overrightarrow{a}$ thì tọa độ $\overrightarrow{a}$ cũng là tọa độ điểm $A$.  Suy ra $A(1;-4;3)$.
			\itemch $\overrightarrow{b}=\overrightarrow{0}\Leftrightarrow\heva{&m-n=0\\&4m-6n=0\\&n^{2}-3m+2=0}\Leftrightarrow\heva{&m=0\\&n=0\\&n^{2}-3m+2=0}\quad\text{ (vô  nghiệm).}$\\
			Vậy không  tồn tại $m,n$ để $\overrightarrow{b}=\overrightarrow{0}$.\\
			\itemch  $\overrightarrow{a}=\overrightarrow{b}\Leftrightarrow\heva{&m-n=1\\&4m-6n=-4\\&n^{2}-3m+2=3}\Leftrightarrow\heva{&m=5\\&n=4.}$\\
			Suy  ra $m+n=9$
		\end{itemchoice}
	 }
\end{ex}

% nguồn: Lop_12_Hinh_hoc_Chuong_2
\begin{ex}%[2H2H2-3]%[ma:f61b7653]
	Trong không gian $Oxyz$, cho $\overrightarrow{a}=(2;0;2)$, $\overrightarrow{b}=2\overrightarrow{j}+2\overrightarrow{k}$. Dựng $\overrightarrow{OA}=\overrightarrow{a}$ và $\overrightarrow{OB}=\overrightarrow{b}$. 
	 \choiceTF
		{$\overrightarrow{a}=2\overrightarrow{i}+2\overrightarrow{k}$} 
		 {\True Toạ độ $\overrightarrow{b}=(0;2;2)$}
		{\True  Toạ độ $\overrightarrow{AB}=(-2;2;0)$}
		{Góc $\widehat{AOB}=45^\circ$} 
	\loigiai{ 
		\begin{center}
			\begin{tikzpicture}[scale=0.8,font=\footnotesize,line join=round,line cap=round,>=stealth]
					\coordinate  (O) at (0,0);
					\coordinate  (H) at (-2,-1);
					\coordinate (K) at  (4,0);
					\coordinate (M) at ($(H)+(K)-(O)$); 
					\coordinate (P) at at  ($(O)+(0,4)$);
					\coordinate  (A) at ($(H)+(0,4)$);
					\coordinate  (N) at ($(M)+(0,4)$);
					\coordinate (B) at  ($(K)+(0,4)$);
					\coordinate (x)  at ($(O)!1.5!(H)$);
					\coordinate (y) at  ($(O)!1.5!(K)$);
					\coordinate (z) at  ($(O)!1.3!(P)$);
					\foreach \i in  {O,H,M,K}{\coordinate (\i') at ($(\i)+(0,4)$);}
					 \draw (P)--(A)--(N)--(B)--cycle;
					\draw [->] (H)--(x); 
					\draw  [->] (K)--(y);
					\draw [->] (P)--(z); 
					\draw (H)--(A) 
						(M)--(N) 
						 (K)--(B)
						(H)--(M)--(K); 
					\draw[dashed,thin](H)--(O)--(P) 
						 (O)--(K);
					\foreach  \i/\g in {P/135,A/180,N/90,B/45,O/-90,H/-90,M/-90,K/-90}{\draw (\i) circle (1pt)
						 ($(\i)+(\g:4mm)$) node[scale=1]{$\i$};}
					\foreach \i/\g in {x/90,y/90,z/0}{\draw ($(\i)+(\g:3mm)$)  node{$\i$};}
				\end{tikzpicture}
		\end{center}
		\begin{itemchoice}
			\itemch Ta có $\overrightarrow{a}=(2;0;2)\Rightarrow  \overrightarrow{a}=2\overrightarrow{i}+2\overrightarrow{k}$.
			\itemch  Ta có $\overrightarrow{b}=2\overrightarrow{j}+2\overrightarrow{k}\Rightarrow \overrightarrow{b}=(0;2;2)$.
			\itemch Ta có $\overrightarrow{OA}=\overrightarrow{a}$ thì toạ độ véc tơ $\overrightarrow{a}$ cũng chính là tọa  độ $A\Rightarrow A(2;0;2)$.\\
			Tương tự $B(0;2;2)$. Từ đây,  ta tính được $\overrightarrow{AB}=(-2;2;0)$.
			\itemch Ta thấy $\widehat{AOB}$ là góc giữa hai véc-tơ $\vec{OA}$  và véc-tơ $\vec{OB}$ hay góc giữa hai véc-tơ $\vec{a}$ và véc-tơ $\vec{b}$.\\
			Ta có $\cos \left(\overrightarrow{a},\overrightarrow{b}\right)=\dfrac{2\cdot  0+0\cdot 2+2\cdot 2}{\sqrt{2^{2}+0^{2}+2^{2}}\cdot \sqrt{0^{2}+2^{2}+2^{2}}}=\dfrac{1}{2}\Rightarrow\left(\overrightarrow{a},\overrightarrow{b}\right)=60^{\circ}$.\\
			Vậy  $\widehat{AOB}=60^{\circ}$.
		\end{itemchoice}
	} 
\end{ex}

% nguồn: Lop_12_Hinh_hoc_Chuong_2
\begin{ex}%[2H2H2-3]%[ma:40240e93]
	Cho hình chóp $S.ABCD$ có đáy $ABCD$ là hình chữ nhật, $AB=1$, $AD=2$, $SA$ vuông góc với mặt đáy và $SA=3$. Với hệ toạ độ $Oxyz$  được thiết lập như sau: Gốc tọa độ $O$ trùng với điểm $A$, các véc tơ $\overrightarrow{AB},\overrightarrow{AD},\overrightarrow{AS}$ lần lượt cùng hướng với $\vec{i},\vec{j}$ và $\vec{k}$. Xét tính đúng sai của các khẳng định sau
	\begin{center}
		\begin{tikzpicture}[scale=0.8,font=\footnotesize,line join=round,line cap=round,>=stealth]
				\coordinate (A) at  (0,0);
				\coordinate (B) at (-2,-2); 
				\coordinate (D) at  (5,0);
				\coordinate (C) at  ($(B)+(D)-(A)$);
				\coordinate (S) at  ($(A)+(0,4)$);
				\coordinate (I) at  ($(A)!0.5!(C)$);
				\foreach  \i in {B,C,D}{\draw (S)--(\i);}
				\draw  (B)--(C)--(D);
				\draw[dashed,thin](S)--(A) 
					(B)--(A)--(D) 
					(A)--(C) 
					 (B)--(D);
				\pic[draw,thin,angle radius=3mm] {right angle  = S--A--D} pic[draw,thin,angle radius=3mm] {right angle = S--A--B};
				\foreach \i/\g in {S/90,A/-90,B/-90,C/-90,D/-90,I/-80}{\draw  (\i) circle (1.5pt)
					($(\i)+(\g:3mm)$) node[scale=1]{$\i$};} 
			\end{tikzpicture}
	\end{center}
	 \choiceTF
		{\True  Tọa độ $D\left(0;2;0\right)$}{Tọa độ $C\left(1;2;3\right)$}
		{\True Tọa độ $S\left(2;0;0\right)$}{Tọa độ $I\left(1;1;0\right)$} 
	\loigiai{ 
		\begin{center}
			\begin{tikzpicture}[scale=0.8,font=\footnotesize,line join=round,line cap=round,>=stealth]
					\coordinate (A)  at (0,0);
					 \coordinate (B) at (-2,-2);
					 \coordinate (D) at (5,0);
					 \coordinate (C) at ($(B)+(D)-(A)$);
					\coordinate (S) at  ($(A)+(0,4)$);
					\coordinate  (I) at ($(A)!0.5!(C)$);
					\foreach  \i in {B,C,D}{\draw (S)--(\i);}
					\draw  (B)--(C)--(D);
					\draw[dashed,thin](S)--(A) 
						 (B)--(A)--(D)
						(A)--(C) 
						 (B)--(D);
					\pic[draw,thin,angle radius=3mm] {right  angle = S--A--D} pic[draw,thin,angle radius=3mm] {right angle = S--A--B};
					\foreach \i/\g in  {S/90,A/-90,B/-90,C/-90,D/-90,I/-80}{\draw (\i) circle (1.5pt)
						($(\i)+(\g:3mm)$) node[scale=1]{$\i$};} 
				\end{tikzpicture}
		\end{center}
		Với hệ trục đã chọn  như hình vẽ thì
		\begin{itemchoice}
			\itemch Điểm $D\in Oy$  và $AD=2$ nên $D\left(0;2;0\right)$.
			\itemch Điểm $C\in \left(Oxy\right)$ và có hình chiếu lên $Ox,Oy$ lần lượt là điểm $B$  và $D$. Do $AB=1$ và $AD=2$ nên $C\left(1;2;0\right)$.
			\itemch  Điểm $S\in Oz$ và $AS=3$ nên $S\left(0;0;3\right)$.
			\itemch Điểm $I\in \left(Oxy\right)$ và và có hình chiếu lên $Ox,Oy$ lần lượt  là trung điểm của $AB$ và $AD$ nên $I\left(0,5;1;0\right)$.
		\end{itemchoice}
	 }
\end{ex}

% nguồn: Lop_12_Hinh_hoc_Chuong_2
\begin{ex}%[2H2H2-3]%[ma:d66a40c3]
	Cho hình lập phương $ABCD.A'B'C'D'$ có cạnh bằng $2$. Với hệ toạ độ $Oxyz$ được thiết lập như hình  bên (gốc tọa độ $O$ trùng với tâm hình vuông $ABCD$), hãy xét tính đúng sai của các khẳng định sau
	 \choiceTF
		{Tọa độ $A\left(-1;0;0\right)$}{\True $\overrightarrow{AC'}=\left(2\sqrt{2};0;2\right)$} 
		{\True Tọa độ $D'\left(0;\sqrt{2};2\right)$}{$\overrightarrow{BD'}=\left(0;0;2\right)$} 
	\loigiai{ 
		\begin{tikzpicture}[scale=0.8,font=\footnotesize,line join=round,line cap=round,>=stealth]
				\coordinate (A) at  (0,0);
				 \coordinate (B) at (-2,-2);
				 \coordinate (D) at (6,0);
				\coordinate (C) at  ($(B)+(D)-(A)$);
				\coordinate (O) at  ($(B)!0.5!(D)$);
				\foreach  \i in {A,B,C,D}{\coordinate (\i') at ($(\i)+(0,4)$);}
				\coordinate  (O') at ($(B')!0.5!(D')$);
				\draw (A')--(B')--(C')--(D')--cycle  (A')--(C')
					 (B')--(D');
				\draw (B)--(B') 
					 (C)--(C')
					 (D)--(D')
					 (B)--(C)--(D);
				\draw[dashed,thin](B)--(A)--(A') 
					(A)--(D) 
					 (A)--(C)
					(B)--(D) 
					 (O)--(O');
				\foreach  \i/\g in {A'/110,B'/90,C'/90,D'/90,A/-90,B/-90,C/-90,D/-90,O/-90}{\draw[fill=blue](\i) circle (1.5pt)
					($(\i)+(\g:5mm)$) node[scale=1]{$\i$};} 
				\coordinate  (Z) at ($(O')+(0,2)$);
				\coordinate (Y) at ($(O)!1.5!(D)$); 
				\coordinate (X) at ($(O)!2!(C)$); 
				 \draw[->] (O')--(Z);
				\draw[->]  (C)--(X);
				 \draw[->] (D)--(Y);
				\draw (X)  node[scale=1,left]{$x$};
				\draw (Y)  node[scale=1,above]{$y$};
				 \draw (Z) node[scale=1,left]{$z$};
				\pic[draw,thin,angle radius=3mm] {right angle  = C--O--D};
			\end{tikzpicture}
		Độ dài $AC=2\sqrt{2}$.  Với hệ trục $Oxyz$ đã chọn như hình vẽ thì
		\begin{itemchoice}
			\itemch Điểm $A\in  Ox$, nằm ngược chiều dương và $OA=\sqrt{2}$ nên $A\left(-\sqrt{2};0;0\right)$.
			\itemch Tọa  độ $C'\left(\sqrt{2};0;2\right)$. Suy ra $\overrightarrow{AC'}=\left(2\sqrt{2};0;2\right)$.
			\itemch Điểm $D'$ có hình chiếu vuông góc xuống ($\left. Oxy\right)$ là điểm $D\left(0;\sqrt{2};0\right)$ và $DD'=2$ nên  $D'\left(0;\sqrt{2};2\right)$.
			\itemch Tọa độ  $B\left(0;-\sqrt{2};0\right),D'\left(0;\sqrt{2};2\right)$. Suy ra $\overrightarrow{BD'}=\left(0;2\sqrt{2};2\right)$.
		\end{itemchoice}
	} 
\end{ex}

% nguồn: Lop_12_Hinh_hoc_Chuong_2
\begin{ex}%[2H2H2-3]%[ma:c2a9cd9b]
	Trong không gian với hệ tọa độ $Oxyz$, cho  hình hộp $ABCD.A'B'C'D'$ biết rằng các điểm có tọa độ $A(2;1;0)$, $C(0;3;0)$, $C'(-1;2;1)$, $D'(0;-2;0)$.
	\choiceTF 
		{Tọa độ các điểm $A'$, $B'$ là $A'(1;0;-1)$,  $B'(0;4;2)$}
		{\True Tọa độ các  điểm $B$, $D$ là $B(1;5;1)$, $D(1;-1;-1)$}
		{\True Tọa  độ vectơ $\overrightarrow{AB}$ là $\overrightarrow{AB}=\overrightarrow{i}+4\overrightarrow{j}+\overrightarrow{k}$}
		{\True Tọa độ vectơ $\overrightarrow{AB}$  là $\overrightarrow{B'D}=\overrightarrow{i}-5\overrightarrow{j}-3\overrightarrow{k}$}
	 \loigiai{
		\begin{itemchoice}
			\itemch Gọi tọa độ điểm $A'$ là  $(x;y;z)\Rightarrow \overrightarrow{A'C'}=(-1-x;2-y;1-z)$.\\
			Khi đó $\overrightarrow{AC}=(-2;2;0)$.\\ 
			Vì  $A'C'CA$ là hình bình hành nên $\overrightarrow{A'C'}=\overrightarrow{AC}$.\\
			Suy ra $\heva{&-1-x=-2\\&2-y=2\\&1-z=0}\Leftrightarrow\heva{&x=1\\&y=0\\&z=1}\Rightarrow  A'(1;0;1)$.\\
			Làm tương tự ta có  $B'(0;4;2)$.
			\itemch  $B(1;5;1)$, $D(1;-1;-1)$.
			 \itemch $\overrightarrow{AB}=(1;4;1)\Rightarrow \overrightarrow{AB}=\overrightarrow{i}+4\overrightarrow{j}+\overrightarrow{k}$.
			\itemch  $\overrightarrow{B'D}=(1;-5;-3)\Rightarrow \overrightarrow{B'D}=\overrightarrow{i}-5\overrightarrow{j}-3\overrightarrow{k}$.
		\end{itemchoice}
	} 
\end{ex}

% nguồn: Lop_12_Hinh_hoc_Chuong_2
\begin{ex}%[2H2H2-3]%[ma:8bf25c10]
	Trong không  gian $Oxyz$, cho ba vectơ $\overrightarrow{a}=(-1;1;0)$, $\overrightarrow{b}=(1;1;0)$, $\overrightarrow{c}=(1;1;1)$.
	\choiceTF 
		{\True  Độ dài vectơ $\overrightarrow{a}$ bằng $\sqrt{2}$}
		{\True  Vectơ $\overrightarrow{b}$ vuông góc với $\overrightarrow{a}$}
		 {Vectơ $\overrightarrow{b}$ vuông góc với $\overrightarrow{c}$}
		{\True Tọa độ vectơ $3\overrightarrow{a}+2\overrightarrow{b}-\overrightarrow{c}$ bằng $(-2;4;-1)$} 
	\loigiai{ 
		\begin{itemchoice}
			\itemch $|\overrightarrow{a}|=\sqrt{(-1)^{2}+1^{2}+0}=\sqrt{2}$. 
			\itemch  $\overrightarrow{a}\cdot\overrightarrow{b}=(-1)\cdot 1+1\cdot 1+0=0$.
			 \itemch $\overrightarrow{b}\cdot\overrightarrow{c}=1\cdot 1+1\cdot 1+0\cdot 1=2$.
			\itemch $3\overrightarrow{a}+2\overrightarrow{b}-\overrightarrow{c}=(-2;4;-1)$. 
		\end{itemchoice}
	} 
\end{ex}

% nguồn: Lop_12_Hinh_hoc_Chuong_2
\begin{ex}%[2H2V2-3]%[ma:4899dd4b]
	Trong không gian $Oxyz$, cho hình hộp $ABCD.A'B'C'D'$ có $A(0;0;0)$, $B(3;0;0)$, $D(0;3;0)$, $D'(0;3;-3)$. Gọi  $G=(a;b;c)$ là trọng tâm tam giác $A'B'C$.
	\choiceTF 
		{$\overrightarrow{BD}=(3;-3;0)$} 
		{\True  $\overrightarrow{AC}=(3;3;0)$}
		{\True  $\overrightarrow{A C'}=(3;3;-3)$}
		{\True $a+2b+c=2$} 
	 \loigiai{
		\begin{itemchoice}
			\itemch $\overrightarrow{BD}=(-3;3;0)$. 
			\itemch $\overrightarrow{AB}=(3;0;0), \overrightarrow{AD}=(0;3;0)\Rightarrow  \overrightarrow{AB}+\overrightarrow{AD}=\overrightarrow{AC}=(3;3;0)$.
			 \itemch $\overrightarrow{AB}=(3;0;0), \overrightarrow{AD}=(0;3;0), \overrightarrow{AA'}=\overrightarrow{DD'}=(0;0;-3) \Rightarrow\overrightarrow{AB}+\overrightarrow{AD}+\overrightarrow{AA'}=\overrightarrow{AC'}=(3;3;-3)$.
			\itemch 
			\begin{center}
				\begin{tikzpicture}[scale=0.75,font=\footnotesize,line join=round,line cap=round]
						\foreach \x\y\t in {0/0/A,-2/-1.5/B,3.9/0/D,-0.5/3.5/A'} 
						\coordinate (\t) at  (\x,\y);
						\coordinate (C) at ($(B)+(D)-(A)$); 
						\coordinate (B')  at ($(A')+(B)-(A)$);
						\coordinate  (C') at ($(B')+(C)-(B)$);
						\coordinate (D') at  ($(C')+(D)-(C)$);
						 \draw (A')--(B')--(B)--(C)--(C');
						\draw (A')--(D')--(D); 
						 \draw (D')--(C')
						(C)--(D); 
						\draw (B')--(C') 
						(D)--(D'); 
						\draw[dashed]  (A)--(B)
						(A')--(A)--(D); 
						\foreach \t/\g in {A/180,B/180,C/0,D/0,A'/180,B'/180,C'/0,D'/0} \draw (\t)  node[shift={(\g:10pt)}]{$\t$};
					\end{tikzpicture}
			\end{center}
			Ta có $\overrightarrow{AB}=(3;0;0)$. Gọi  $C(x;y;z) \Rightarrow \overrightarrow{DC}=(x;y-3;z)$.\\
			$ABCD$ là hình bình hành $\Leftrightarrow \overrightarrow{AB}=\overrightarrow{DC} \Leftrightarrow\left\{\begin{aligned}&x=3\\&y=3\\ &z=0
		\end{aligned}\right.$. Suy ra $C(3;3;0)$.\\
		Ta lại có $\overrightarrow{AD}=(0;3;0)$.  Gọi $A'\left(x';y';z'\right) \Rightarrow \overrightarrow{A'D'}=\left(-x';3-y';-3-z'\right)$.\\
		$ADD'A'$ là hình bình hành $\Leftrightarrow \overrightarrow{AD}=\overrightarrow{A'D'} \Leftrightarrow\left\{\begin{aligned}&x'=0\\&y'=0\\& z'=-3.
	\end{aligned}\right.$ Suy ra $A'(0;0;-3)$.\\
	Gọi $B'\left(x_0;y_0;z_0\right)\Rightarrow\overrightarrow{A'B'}=\left(x_0;y_0;z_0+3\right)$.\\ 
	$ABB'A'$ là hình bình hành $\Leftrightarrow\overrightarrow{AB}=\overrightarrow{A'B'} \Leftrightarrow\left\{\begin{aligned}&x_0=3\\&y_0=0\\&z_0=-3.
	\end{aligned}\right.$ Suy ra $B'(3;0;-3)$. \\
	$G\left(x_G;y_G;z_G\right)$ là trọng tâm tam giác $A'B'C\Leftrightarrow\left\{\begin{aligned}&x_G=\dfrac{0+3+3}{3}=2\\&y_G=\frac{0+0+3}{3}=1\\&z_G=\dfrac{-3-3+0}{3}=-2.
	\end{aligned}\right.$\\ Suy ra $G(2;1;-2)$.\\
	 Vậy $S=a+2b+c=2+2\cdot1-2=2$.
	\end{itemchoice}
	 }
	\end{ex}

% nguồn: Lop_12_Hinh_hoc_Chuong_2
\begin{ex}%[2H2V2-3]%[ma:6192bc07]
	Trong không gian $Oxyz$, cho các điểm  $A(4; 2;-1), B(1;-1; 2)$ và $C(0;-2; 3)$. Khi đó
	 \choiceTF
		{\True $\overrightarrow{AB}=(-3;-3;  3)$}
		{$\left| \overrightarrow{AB}\right|  =2\sqrt{3}$}
		{\True  Điểm $M$ thỏa mãn $\overrightarrow{AB}+\overrightarrow{CM}=\overrightarrow0$ lúc đó $M(3; 1; 0)$}
		{\True Điểm $N$ thuộc mặt phẳng $(Oxy)$, sao cho $A, B, N$ thẳng hàng thì $N(3; 1; 0)$} 
	\loigiai{ 
		\begin{itemchoice}
			 \itemch Đúng. $
			\overrightarrow{A B}=(1-4;-1-2; 2+1)=(-3;-3; 3)$. 
			\itemch Sai.  $|\overrightarrow{AB}|=2\sqrt{3}$
			$\Rightarrow|\overrightarrow{AB}|=\sqrt{(-3)^2+(-3)^2+3^2}=3\sqrt{3}$. 
			\itemch Đúng. Gọi $M\ (x;  y; z)$ thì $\overrightarrow{MC}=(-x;-2-y, 3-z)$.\\
			Vì  $\overrightarrow{AB}+\overrightarrow{CM}=\overrightarrow0\Rightarrow \overrightarrow{AB}=\overrightarrow{MC} \Rightarrow\heva{&-x=-3\\
			 &-2-y=-3\\
			 &3-z=3} \Leftrightarrow\heva{&x=3\\
			&y=1. \\ 
			&z=0} \Rightarrow M(3; 1; 0)$. 
			\itemch Đúng. Vì $N$ thuộc mặt phẳng $(Oxy)$ nên tọa độ điểm $N$ là $N(x;  y; 0)$.\\
			Ta có $\overrightarrow{AN}(x-4; y-2; 1); \overrightarrow{BN}(x-1; y+1;-2)$.\\ 
			Để $A$,  $B$, $N$ thẳng hàng thì hai vectơ $\overrightarrow{AN}$, $\overrightarrow{BN}$ cùng phương. Do đó, $\overrightarrow{AN}=k \overrightarrow{BN}$\\
			$$\Rightarrow\heva{&x-4=k(x-1) \\ 
			 &y-2=k(y+1) \\
			&1=-2k}  \Rightarrow\heva{&x-4=-\dfrac{1}{2}(x-1) \\
			&y-2=-\dfrac{1}{2}(y+1)  \\
			&k=\dfrac{-1}{2}}  \Rightarrow\heva{&x=3\\
			&y=1.}$$ 
			Vậy $N(3; 1;  0)$.
		\end{itemchoice}
	 }
\end{ex}

% nguồn: Lop_12_Hinh_hoc_Chuong_2
\begin{ex}%[2H2V2-3]%[ma:e67f545f]
	Cho  tứ diện $ABCD$ với $A\left(2;1;0\right)$, $B\left(1;1;3\right)$, $C\left(2;-1;3\right)$, $D\left(1;-1;0\right)$.
	 \choiceTF
		{\True Tứ diện $ABCD$ có  các cạnh đối đôi một bằng nhau}
		{Góc giữa $2$ đường thẳng $AB$  và $CD$ là $\varphi=\arccos (0{,}3)$}
		{Khoảng cách giữa $2$ đường thẳng $AB$ và $CD$ bằng  $3$}
		{\True Bán kính của mặt cầu ngoại tiếp tứ diện $ABCD$ bằng  $\dfrac{\sqrt{14}}{2}$}
	\loigiai{ 
		\begin{itemchoice}
			\itemch Ta  có $AB=CD=\sqrt{10}$; $AC=BD=\sqrt{13}$; $AD=BC=\sqrt{5}$.\\
			Vậy tứ diện $ABCD$ có  các cạnh đối đôi một bằng nhau.
			\itemch Ta có $\overrightarrow{AB}=\left(-1;0;3\right)$, $\overrightarrow{CD}=\left(-1;0;-3\right)$.  Gọi $\varphi$ là góc giữa $AB$ và $CD$.
			\[\cos  \varphi=\cos \left(\overrightarrow{AB},\overrightarrow{CD}\right)=\dfrac{\left|-8\right|}{\sqrt{10}\cdot\sqrt{10}}=\dfrac{4}{5}.
			\] 
			Vậy  góc giữa $AB$ và $CD$ là $\varphi=\arccos 0{,}8$.
			\itemch Lấy $I$ trung điểm của $AB$, $J$  là trung điểm của $CD$.\\
			Ta có $\triangle ACD=\triangle BCD$ (c-c-c)  nên $2$ đường trung tuyến tương ứng $AJ=BJ$.\\
			 Vậy $\triangle AJB$ cân đỉnh $J$ nên $IJ$ vuông góc với $AB$ tại $I$.\\
			 Tương tự $\triangle ICD$ cân đỉnh $I$ nên $IJ$ vuông góc với $CD$ tại $J$.\\
			Vậy $IJ$ là đường  vuông góc chung của $AB$ và $CD$ ta được $I\left(\dfrac{3}{2};1;\dfrac{3}{2} \right)$ và $J\left(\dfrac{3}{2};-1;\dfrac{3}{2} \right)$.\\
			Vậy khoảng cách giữa $AB$ và $CD$  chính là độ dài đoạn vuông góc chung $IJ$.
			\[d\left(AB;CD\right)=IJ=\sqrt{\left(\dfrac{3}{2}-\dfrac{3}{2}  \right)^2+\left[1-\left(-1\right)\right]^2+\left(\dfrac{3}{2}-\dfrac{3}{2} \right)^2}=2.\]
			\itemch Theo kết  quả câu c). Lấy $G$ là trung điểm của $IJ$ ta được\\
			$GA=GB$ vì $\triangle GAB$ cân đỉnh $G$; $GC=GD$  vì $\triangle GCD$ cân đỉnh $G$.\\
			Mà  $GA=\sqrt{GI^2+IA^2}$ mà $GI=GJ,IA=ID$ và $GC=\sqrt{GJ^2+ID^2}$.\\
			Do  đó $GA=GB=GC=GD=R$.\\
			 Do đó $G$ tâm mặt cầu ngoại tiếp khối tứ diện $ABCD$ và $G\left(\dfrac{3}{2};0;\dfrac{3}{2} \right)$ và bán kính của mặt cầu là $R=GA=\dfrac{\sqrt{14}}{2}$ ($G$ cũng chính là trọng tâm của khối tứ diện gần đều $ABCD$).
		\end{itemchoice}
	} 
\end{ex}

% nguồn: Lop_12_Hinh_hoc_Chuong_2
\begin{ex}%[2H2V2-3]%[ma:ffffaf3a]
	Trong không gian $Oxyz$, cho tam giác $ABC$ có $A\left( 1; -1; 0 \right)$, $B\left( -2; 5;  3 \right)$, $C\left( 3; 4; 9 \right)$.
	 \choiceTF
		{Toạ độ trọng tâm $G$  của tam giác $ABC$ là $G\left( \dfrac{2}{3}; \dfrac{4}{3};4 \right)$}
		{Tọa độ vectơ $\overrightarrow{AB}=\left( 3;  -6; -3 \right)$}
		{\True Với  điểm $D\left( 6; -2; 6 \right)$ thì tứ giác $ABCD$ là hình bình hành}
		{Xét điểm $M\left( a; b; c \right)$ thuộc  đoạn $AB$ sao cho $MA=2MB$. Khi đó $a+b+c=6$}
	\loigiai{ 
		\begin{itemchoice}
			 \itemch {\bf Sai.}\\
			Toạ độ trọng  tâm $G$ của tam giác $ABC$ là $G\left( \dfrac{2}{3}; \dfrac{8}{3};4 \right)$.
			\itemch {\bf Sai.}\\ 
			Tọa độ vectơ $\overrightarrow{AB}=\left( -3;  6; 3 \right)$.
			\itemch  {\bf Đúng.}\\
			Tứ  giác $ABCD$ là hình bình hành khi $D(6; -2; 6)$
			\itemch {\bf  Sai.}\\
			Điểm $M$ thuộc  đoạn $AB$ sao cho $MA=2MB$.\\
			Khi đó  $\overrightarrow{AM}=\dfrac{2}{3}\overrightarrow{AB}=(-2; 4; 2)$. Điểm $A(1; -1; 0)$ nên $M\left( -1; 3; 2 \right)$.\\
			Vậy  $a+b+c=4$.
		\end{itemchoice}
	} 
\end{ex}


%%%%%%%%%%%%%%%% TRẢ LỜI NGẮN %%%%%%%%%%%%%%%%

% nguồn: Lop_12_Hinh_hoc_Chuong_2
\begin{bt}%[2H2N2-3]%[ma:6bd4b674]
	Trong không gian $Oxyz$, cho hai điểm $A(1;-2;5)$ và $B(-2;-2;1)$. Tính  độ dài đoạn thẳng $AB$.
	 \shortans{5}
	\loigiai{ 
		 Ta có $A B=\sqrt{(-2-1)^2+(-2+2)^2+(1-5)^2}=5$.
	 }
\end{bt}

% nguồn: Lop_12_Hinh_hoc_Chuong_2
\begin{bt}%[2H2H2-3]%[ma:3d684955]
	Trong không gian $Oxyz$, cho hình hộp $ABCD.A'B'C'D'$ có tọa độ các điểm $A\left(1;2;-1\right),C\left(3;-4;1\right),B'\left(2;-1;3\right),D'\left(0;3;5\right)$. Giả sử tọa  độ điểm $A'\left(x;y;z\right)$ thì $x+y+z$ là
	 \shortans{7}
	 \loigiai{
		\begin{center}
			\begin{tikzpicture}[scale=0.6,font=\footnotesize,line join=round,line cap=round,>=stealth]
					 \path
						(0,0)  coordinate (A)
						++(-130:3)  coordinate (B)
						 ++(0:4) coordinate (C)
						 ($(A)+(C)-(B)$) coordinate (D)
						($(A)!1/2!(C)$) coordinate (O); 
					 \foreach \i in {A,B,C,D}{
						\coordinate  (\i') at ($(\i)+(1,4)$);
					 }
					 \draw (A')--(B')--(C')--(D')--cycle;
					\draw  (B)--(B')
						 (C)--(C')
						 (D)--(D')
						 (B)--(C)--(D)
						(A')--(C'); 
					\draw[dashed,thin](B)--(A)--(A') 
						(C)--(A)--(D)--(B); 
					\foreach  \i/\g in {A'/90,B'/90,C'/90,D'/90,A/-90,B/-90,C/-90,D/-90}
					\fill (\i)  circle (1pt)+(\g:5mm) node[scale=1]{$\i$};
				\end{tikzpicture}
		\end{center}
		Theo quy  tắc hình hộp, ta có
		\begin{eqnarray*}
			&&\overrightarrow{A'B'}+\overrightarrow{A'D'}+\overrightarrow{A'A}=\overrightarrow{A'C}\\ 
			&\Leftrightarrow&  \heva{&\left(2-x\right)+\left(0-x\right)+\left(1-x\right)=3-x \\&\left(-1-y\right)+\left(3-y\right)+\left(2-y\right)=-4-y \\&\left(3-z\right)+\left(5-z\right)+\left(-1-z\right)=1-z}\\
			&\Leftrightarrow&\heva{&x=0 \\&y=4  \\&z=3.}
		\end{eqnarray*}
	Suy  ra $x+y+z=0+4+3=7$.}
\end{bt}

% nguồn: Lop_12_Hinh_hoc_Chuong_2
\begin{bt}%[2H2H2-3]%[ma:16e3a9eb]
	Trong không gian $Oxyz$, cho hình thang $ABCD$ vuông tại $A$ và $B$. Ba đỉnh $A(1;2;1)$, $B(2;0;-1)$, $C(6;1;0)$ và hình thang có diện tích bằng $6\sqrt{2}$.  Giả sử đỉnh $D(a;b;c)$. Tính $a+b+c$.
	\shortans{6} 
	 \loigiai{
		\begin{center}
			\begin{tikzpicture}[scale=1,font=\footnotesize,line join=round,
						line  cap=round,>=stealth]
					 \draw (0:0) coordinate (C)--++(180:4.5) coordinate (B)--++(90:2.5) coordinate (A)--++(0:1.5) coordinate (D)--cycle;
					\foreach\x/\y/\z in {D/A/B,A/B/C}\draw pic[draw,angle radius=2mm] 
						{right angle=\x--\y--\z}; 
					\foreach  \x/\g in {A/150,B/180,C/0,D/30}
					 \fill (\x) circle (1pt)
						($(\x)+(\g:3mm)$) node{$\x$}; 
				\end{tikzpicture}
		\end{center}
		Ta  có $\overrightarrow{AB}=(1;-2;-2) \Rightarrow\left| \overrightarrow{AB}\right| =3$; $\overrightarrow{BC}=(4;1;1) \Rightarrow\left| \overrightarrow{BC}\right| =3 \sqrt{2}$.\\
		Theo giả thiết $ABCD$ là hình thang vuông tại $A$ và $B$ và có diện tích bằng $6\sqrt{2}$ nên $$\dfrac{1}{2}\cdot AB\cdot (AD+BC)=6\sqrt{2}  \Leftrightarrow \dfrac{1}{2} \cdot 3 \cdot(AD+3\sqrt{2})=6\sqrt{2} \Rightarrow AD=\sqrt{2} \Rightarrow AD=\dfrac{1}{3} BC.$$
		Do $ABCD$ là hình thang vuông tại $A$ và  $B$ nên $\overrightarrow{AD}=\dfrac{1}{3} \overrightarrow{BC}$.\\
		Giả sử $D(a;b;c)$. Khi đó ta có $\heva{&a-1=\dfrac{4}{3} \\& b-2=\dfrac{1}{3} \\ &c-1=\dfrac{1}{3}} \Leftrightarrow\heva{&a=\dfrac{7}{3}  \\& b=\dfrac{7}{3} \\& c=\dfrac{4}{3}} \Rightarrow a+b+c=6$.
	 }
\end{bt}

% nguồn: Lop_12_Hinh_hoc_Chuong_2
\begin{bt}%[2H2H2-3]%[ma:9c609d5f]
	Cho  tứ diện $O.ABC$ có $OA$, $OB$, $OC$ đôi một vuông góc với nhau $OA=a$ và $OB=OC=2a$. Gọi $P$ là trung điểm của $BC$ (minh họa như hình vẽ). Khi $a=1$ xác định tung độ của vectơ $\overrightarrow{AB}$
	\begin{center}
		\begin{tikzpicture}[scale=0.8,font=\footnotesize,line join=round,line cap=round,>=stealth]
				 \coordinate (O) at (0,0);
				\coordinate  (C) at (-2,-2);
				\coordinate (B)  at (6,0);
				\coordinate (P) at  (2,-1);
				\coordinate  (A) at ($(O)+(0,4)$);
				\draw[dashed]  (A)--(O)--(C)
					 (B)--(O)--(P);
				\draw  (C)--(A)--(B)--(C);
				\foreach \i/\g in {C/110,A/180,B/90,O/180,P/0}{\draw[fill=blue](\i) circle  (1.5pt)
					 ($(\i)+(\g:5mm)$) node[scale=1]{$\i$};}
				 \coordinate (Z) at ($(A)+(0,1)$);
				\coordinate (Y) at ($(B)+(1,0)$); 
				 \coordinate (X) at ($(C)-(1,1)$);
				\draw[->]  (A)--(Z);
				\draw[->]  (C)--(X);
				\draw[->] (B)--(Y); 
				\draw (X) node[scale=1,left]{$x$}; 
				\draw (Y) node[scale=1,above]{$y$}; 
				\draw  (Z) node[scale=1,left]{$z$};
				\draw  ($(O)-(1,1)$) node[scale=1,left]{$2a$};
				\draw ($(O)+(0,2)$)  node[scale=1,right]{$a$};
				\draw  ($(O)+(3,0)$) node[scale=1,above]{$2a$};
			\end{tikzpicture}
	\end{center}
	 \shortans{2}
	 \loigiai{
		\begin{center}
			\begin{tikzpicture}[scale=0.8,font=\footnotesize,line join=round,line cap=round,>=stealth]
					 \coordinate (O) at (0,0);
					\coordinate (C)  at (-2,-2);
					\coordinate (B)  at (6,0);
					\coordinate (P) at (2,-1); 
					 \coordinate (A) at ($(O)+(0,4)$);
					\draw[dashed] (A)--(O)--(C) 
						 (B)--(O)--(P);
					\draw (C)--(A)--(B)--(C); 
					\foreach \i/\g in {C/110,A/180,B/90,O/180,P/0}{\draw[fill=blue](\i) circle (1.5pt) 
						($(\i)+(\g:5mm)$) node[scale=1]{$\i$};} 
					\coordinate (Z)  at ($(A)+(0,1)$);
					\coordinate (Y) at  ($(B)+(1,0)$);
					\coordinate (X)  at ($(C)-(1,1)$);
					\draw[->]  (A)--(Z);
					\draw[->] (C)--(X); 
					 \draw[->] (B)--(Y);
					\draw (X)  node[scale=1,left]{$x$};
					\draw (Y) node[scale=1,above]{$y$}; 
					\draw  (Z) node[scale=1,left]{$z$};
					\draw ($(O)-(1,1)$) node[scale=1,left]{$2a$}; 
					 \draw ($(O)+(0,2)$) node[scale=1,right]{$a$};
					\draw  ($(O)+(3,0)$) node[scale=1,above]{$2a$};
				\end{tikzpicture}
		\end{center}
		\noindent Đặt khối  tứ diện $O.ABC$ vào hệ trục tọa độ gốc $Oxyz$ các điểm $A$, $B$, $C$ lần lượt thuộc các trục tọa độ $Oz,Oy,Ox$.\\
		Ta có $O\left(0;0;0\right)$, $A\left(0;0;1\right)$,  $B\left(0;2;0\right)$, $C\left(2;0;0\right)$, $P\left(1;1;0\right)$, $\overrightarrow{AB}=\left(0;2;-1\right)$.
	 }
\end{bt}

% nguồn: Lop_12_Hinh_hoc_Chuong_2
\begin{bt}%[2H2H2-3]%[ma:d790b906]
	Trong không gian $Oxyz$, cho ba điểm $A(3;2;-1)$, $B(-1;-x;1)$, $C(7;-1;y)$.  Khi $A$, $B$, $C$ thẳng hàng thì giá trị biểu thức $x+y$ bằng bao nhiêu?
	 \shortans{-8}
	 \loigiai{
		 Ta có $\overrightarrow{AB}=(-4;-x-2;2)$, $\overrightarrow{AC}=(4;-3;y+1)$.\\
		Để $A$, $B$, $C$ thẳng hàng  thì $\overrightarrow{AB}=k \overrightarrow{AC} \Leftrightarrow\heva{&-4=k \cdot 4 \\ &-x-2=k \cdot(-3) \\ & 2=k \cdot(y+1)} \Leftrightarrow
		\heva{&k=-1 \\& x=-5 \\& y=-3.}$\\ 
		Vậy  $x+y=-5-3=-8$.
	} 
\end{bt}

% nguồn: Lop_12_Hinh_hoc_Chuong_2
\begin{bt}%[2H2H2-3]%[ma:69415cfc]
	Trong không gian $Oxyz$, cho hai điểm $A(1;-2;3)$, $B(4;1;-1)$. Điểm $M(a;b;c)$ thỏa mãn $MA \cdot \overrightarrow{MA}=4MB\cdot \overrightarrow{MB}$. Giá tri biểu thức $a+b+c$  bằng bao nhiêu?
	 \shortans{6}
	\loigiai{ 
		Vì $MA\cdot \overrightarrow{MA}=4MB\cdot \overrightarrow{MB}$ nên $\overrightarrow{MA}$ cùng hướng $\overrightarrow{MB}$.\\ 
		Vì $MA \cdot\overrightarrow{MA}=4MB\cdot  \overrightarrow{MB}$ nên $MA^2=4MB^2\Leftrightarrow MA=2MB$.\\
		Từ đó có $B$  là trung điểm $AM$. Suy ra $M(7;4;-5)$.\\
		 Vậy $a+b+c=7+4+(-5)=6$.
	} 
\end{bt}

% nguồn: Lop_12_Hinh_hoc_Chuong_2
\begin{bt}%[2H2H2-3]%[ma:8c8bafee]
	Trong không gian $Oxy$, cho ba điểm $A(1;-1;1)$, $B(3;1;2)$ và $C(-1;0;3)$. Có bao nhiêu điểm $D$ sao cho tứ giác $ABCD$ là hình thang có $2$ cạnh đáy $AB$, $CD$ và có góc tại  $D$ bằng $45^{\circ}$?
	 \shortans{2}
	 \loigiai{
		\begin{center}
			\begin{tikzpicture}[scale=1,font=\footnotesize,line join=round,
						line  cap=round,>=stealth]
					\draw (0:0)  coordinate (D)--++(45:3) coordinate (A)--++(0:1) coordinate (B)
						(D)--++(0:4) coordinate (C)--(B); 
					\draw pic[draw,angle radius=5mm]{angle= C--D--A}; 
					\draw ($(D)+(22:.8)$) node{$45^\circ$}; 
					\foreach  \x/\g in {A/150,B/30,C/0,D/-180}
					 \fill (\x) circle (1pt)
						 ($(\x)+(\g:3mm)$) node{$\x$};
				\end{tikzpicture}
		\end{center}
		Gọi  $D(x;y;z)$. Ta có $\overrightarrow{AB}=(2;2;1)$, $\overrightarrow{DC}=(-1-x;-y;3-z)$.\\
		Vì tứ giác $ABCD$ là hình thang có $2$ cạnh đáy $AB$, $CD$ nên $\overrightarrow{AB}$, $\overrightarrow{DC}$  cùng phương.\\
		Từ đó có $\dfrac{-1-x}{2}=\dfrac{-y}{2}=\dfrac{3-z}{1}  \Rightarrow\heva{&x=2z-7 \\&y=2z-6.}$\\
		Lúc đó $\overrightarrow{AD}=(x-1;y+1;z-1)=(2z-8;2z-5;z-1)$ và $\overrightarrow{CD}=(x+1;y;z-3)=(2z-6;2z-6;z-3)$.\\ 
		Ta có $\cos \left(\overrightarrow{AD},\overrightarrow{CD}\right)=\dfrac{9(z-3)^2}{\sqrt{9\left(z^2-6z+10\right)} \cdot  \sqrt{9(z-3)^2}}=\dfrac{z-3}{\sqrt{z^2-6 z+10}}=\dfrac{1}{\sqrt{2}}$.\\
		Suy ra $2\left(z^2-6z+9\right)=z^2-6z+10\Leftrightarrow\hoac{&z=4\\&z=2}\Rightarrow \heva{&x=1;y=2 \\&x=-3;y=-2}$.\\ 
		 Vậy $D(1;2;4)$ hoặc $D(-3;-2;2)$.
	} 
\end{bt}

% nguồn: Lop_12_Hinh_hoc_Chuong_2
\begin{bt}%[2H2H2-3]%[ma:0ca349a8]
	Cho tứ diện $SABC$ có $SC=CA=AB=3\sqrt{2}$, $SC$ vuông góc $\left(ABC\right)$, tam giác  $ABC$ vuông tại $A$, các điểm $M$ và $N$ lần lượt thuộc $SA$ và $BC$ sao cho $AM=CN=2$. Tung độ của $\overrightarrow{NB}$ khi đó bằng bao nhiêu?\\
	Chọn hệ tọa độ như hình vẽ. 
	\begin{center}
		\begin{tikzpicture}[scale=0.8,font=\footnotesize,line join=round,line cap=round,>=stealth]
				\coordinate (C) at  (0,0);
				\coordinate  (A') at (-2,-2);
				\coordinate (B) at  (6,0);
				\coordinate (B') at  (4,0);
				\coordinate  (A) at ($(B')+(A')-(C)$);
				\coordinate (S)  at ($(C)+(0,4)$);
				\coordinate  (M) at ($(S)!0.6!(A)$);
				\coordinate (N) at  ($(C)!0.4!(B)$);
				 \draw[dashed] (A')--(A)--(B');
				 \draw (C)--(A)--(S)--(B)--(A);
				\draw[dashed,red] (M)--(N); 
				 \foreach \i/\g in {C/110,A/-90,B/90,S/180,M/180,N/-90}{\draw[fill=blue](\i) circle (1.5pt)
					($(\i)+(\g:5mm)$)  node[scale=1]{$\i$};}
				 \coordinate (Z) at ($(C)+(0,5)$);
				\coordinate (Y) at  ($(B)+(1,0)$);
				\coordinate  (X) at ($(A')-(1,1)$);
				\draw[->]  (C)--(Z);
				\draw[->] (C)--(X); 
				\draw[->] (C)--(Y); 
				\draw (X)  node[scale=1,left]{$x$};
				\draw (Y)  node[scale=1,above]{$y$};
				\draw  (Z) node[scale=1,left]{$z$};
			\end{tikzpicture}
	\end{center}
	 \shortans{4}
	\loigiai{ 
		\begin{center}
			\begin{tikzpicture}[scale=0.8,font=\footnotesize,line join=round,line cap=round,>=stealth]
					\coordinate (C) at  (0,0);
					\coordinate (A') at  (-2,-2);
					 \coordinate (B) at (6,0);
					\coordinate (B') at (4,0); 
					\coordinate (A) at  ($(B')+(A')-(C)$);
					\coordinate (S) at  ($(C)+(0,4)$);
					\coordinate (M) at  ($(S)!0.6!(A)$);
					\coordinate (N)  at ($(C)!0.4!(B)$);
					\draw[dashed] (A')--(A)--(B'); 
					 \draw (C)--(A)--(S)--(B)--(A);
					\draw[dashed,red]  (M)--(N);
					\foreach \i/\g in {C/110,A/-90,B/90,S/180,M/180,N/-90}{\draw[fill=blue](\i)  circle (1.5pt)
						($(\i)+(\g:5mm)$) node[scale=1]{$\i$};} 
					\coordinate  (Z) at ($(C)+(0,5)$);
					\coordinate  (Y) at ($(B)+(1,0)$);
					\coordinate (X)  at ($(A')-(1,1)$);
					\draw[->] (C)--(Z); 
					 \draw[->] (C)--(X);
					\draw[->]  (C)--(Y);
					 \draw (X) node[scale=1,left]{$x$};
					 \draw (Y) node[scale=1,above]{$y$};
					\draw (Z) node[scale=1,left]{$z$}; 
				\end{tikzpicture}
		\end{center}
		Chọn hệ trục tọa độ $Oxyz$ thỏa $C\left(0;0;0\right)\equiv O$ như hình  vẽ.\\
		Điểm  $B\left(0;6;0\right)\in Oy$, $S\left(0;0;3\sqrt{2} \right)\in Oz$, $A\left(3;3;0\right)$.\\
		 Khi đó tọa độ $N\left(0;2;0\right)$; $M\left(2;2;\sqrt{2} \right)$.\\
		 $\overrightarrow{NM}=\left(2;0;\sqrt{2} \right)$; $\overrightarrow{SB}=\left(0;6;-3\sqrt{2} \right)$ và $\overrightarrow{NB}=\left(0;4;0\right)$.
	} 
\end{bt}

% nguồn: Lop_12_Hinh_hoc_Chuong_2
\begin{bt}%[2H2H2-3]%[ma:c5da69bf]
	Cho tứ diện $OABC$ có $OA$, $OB$, $OC$ đôi một vuông góc và $OA=OB=2a$, $OC=a\sqrt{2}$. Khi đó vectơ $\overrightarrow{AB}\left(m; n;  p\right)$. Khi $a=1$, hãy tính giá trị biểu thức $T=m+n+p$.
	\begin{center}
		\begin{tikzpicture}[scale=0.8,font=\footnotesize,line join=round,line cap=round,>=stealth]
				\coordinate (O)  at (0,0);
				\coordinate  (B) at (-2,-2);
				\coordinate  (C) at (6,0);
				\coordinate (A)  at ($(O)+(0,4)$);
				\coordinate (M) at  ($(A)!0.5!(B)$);
				\coordinate (H)  at ($(M)!0.5!(C)$);
				\draw[dashed] (A)--(O)--(C) 
					 (B)--(O)--(H)
					(O)--(M); 
				 \draw (C)--(A)--(B)--(C)--(M);
				 \foreach \i/\g in {C/-90,A/0,B/-30,O/180,H/90}{\draw[fill=blue](\i) circle (1.5pt)
					($(\i)+(\g:5mm)$)  node[scale=1]{$\i$};}
				\coordinate (Z)  at ($(A)+(0,1)$);
				\coordinate (Y) at ($(O)!1.2!(C)$); 
				\coordinate (X)  at ($(O)!1.2!(B)$);
				\draw[->] (A)--(Z); 
				 \draw[->] (B)--(X);
				\draw[->] (C)--(Y); 
				 \draw (X) node[scale=1,left]{$x$};
				 \draw (Y) node[scale=1,above]{$y$};
				\draw (Z)  node[scale=1,left]{$z$};
				\draw (A) node[scale=1,left]{$2a$}; 
				\draw (B) node[scale=1,left]{$2b$}; 
				 \draw (C) node[scale=1,above]{$a\sqrt{2}$};
			\end{tikzpicture}
	\end{center}
	 \shortans{0}
	\loigiai{ 
		\begin{center}
			\begin{tikzpicture}[scale=0.8,font=\footnotesize,line join=round,line cap=round,>=stealth]
					\coordinate (O) at  (0,0);
					\coordinate (B)  at (-2,-2);
					\coordinate (C)  at (6,0);
					\coordinate (A) at  ($(O)+(0,4)$);
					\coordinate (M) at  ($(A)!0.5!(B)$);
					\coordinate  (H) at ($(M)!0.5!(C)$);
					\draw[dashed] (A)--(O)--(C) 
						 (B)--(O)--(H)
						(O)--(M); 
					\draw  (C)--(A)--(B)--(C)--(M);
					 \foreach \i/\g in {C/-90,A/0,B/-30,O/180,H/90}{\draw[fill=blue](\i) circle (1.5pt)
						 ($(\i)+(\g:5mm)$) node[scale=1]{$\i$};}
					\coordinate (Z) at  ($(A)+(0,1)$);
					\coordinate  (Y) at ($(O)!1.2!(C)$);
					\coordinate  (X) at ($(O)!1.2!(B)$);
					\draw[->] (A)--(Z); 
					\draw[->]  (B)--(X);
					 \draw[->] (C)--(Y);
					\draw (X)  node[scale=1,left]{$x$};
					\draw (Y)  node[scale=1,above]{$y$};
					\draw (Z)  node[scale=1,left]{$z$};
					\draw  (A) node[scale=1,left]{$2a$};
					 \draw (B) node[scale=1,left]{$2b$};
					 \draw (C) node[scale=1,above]{$a\sqrt{2}$};
				\end{tikzpicture}
		\end{center}
		Xét hệ trục tọa độ $Oxyz$ như  sau điểm $O$ là gốc tọa độ $OA\equiv Oz$; $OB\equiv Ox$ và $OC\equiv Oy$.
		Khi đó ta có $O\left(0; 0; 0\right)$; $A\left(0; 0; 2a\right)$;  $B\left(2a; 0; 0\right)$ và $C\left(0; a\sqrt{2}; 0\right)$.
		Vậy ta  có vectơ khi đó là $\overrightarrow{AB}=\left(2a; 0;-2a\right)$ nên $\heva{m&=2a \\ n&=0 \\ p&=-2a} \Rightarrow T=2a+0+\left(-2a\right)=0$.
	} 
\end{bt}

% nguồn: Lop_12_Hinh_hoc_Chuong_2
\begin{bt}%[2H2H2-3]%[ma:7acf23be]
	Trong không gian với hệ tọa độ $Oxyz$ cho các véc tơ $\overrightarrow{u}=2\overrightarrow{i}-2\overrightarrow{j}+\overrightarrow{k}$, $\overrightarrow{v}=(m;2;m+1)$ với $m$ là tham số thực. Có bao nhiêu  giá trị của $m$ để $|\overrightarrow{u}|=|\overrightarrow{v}|$.
	\shortans{2} 
	 \loigiai{
		Ta  có $\overrightarrow{u}=(2;-2;1)$.\\
		Khi đó $|\overrightarrow{u}|=\sqrt{2^{2}+(-2)^{2}+1^{2}}=3$ và  $|\overrightarrow{v}|=\sqrt{m^{2}+2^{2}+(m+1)^{2}}=\sqrt{2m^{2}+2m+5}$.\\
		Do đó $|\overrightarrow{u}|=|\overrightarrow{v}|\Leftrightarrow 9=2m^{2}+2m+5\Leftrightarrow m^{2}+m-2=0\Leftrightarrow\hoac{&m=1\\&m=-2.}$ nên có  hai giá trị.
	 }
\end{bt}

% nguồn: Lop_12_Hinh_hoc_Chuong_2
\begin{bt}%[2H2H2-3]%[ma:3c7124c1]
	Cho hình lập phương $ABCD.A'B'C'D'$ có độ dài cạnh bằng $1$. Gọi $M$, $N$, $P$, $Q$ lần lượt là trung điểm của các cạnh $AB$, $BC$, $C'D'$ và $DD'$. Xác định hoành  độ của vectơ $\overrightarrow{MQ}$. Chọn hệ tọa độ như hình vẽ.
	\begin{center}
		\begin{tikzpicture}[scale=0.8,font=\footnotesize,line join=round,line cap=round,>=stealth]
				\coordinate (D)  at (0,0);
				\coordinate (A) at (-2,-2); 
				\coordinate (C) at  (6,0);
				 \coordinate (B) at ($(A)+(C)-(D)$);
				\foreach  \i in {A,B,C,D}{\coordinate (\i') at ($(\i)+(0,4)$);}
				\draw (A')--(B')--(C')--(D')--cycle; 
				 \draw (A)--(A')
					 (C)--(C')
					 (B)--(B')
					 (A)--(B)--(C);
				\draw[dashed,thin](A)--(D)--(D') 
					(C)--(D); 
				 \coordinate (P) at ($(D')!0.5!(C')$);
				\coordinate (Q) at  ($(D')!0.5!(D)$);
				 \coordinate (M) at ($(A)!0.5!(B)$);
				\coordinate (N) at ($(B)!0.5!(C)$); 
				\draw[dashed,thin](P)--(Q)--(M)--(N)--(P)--(M) 
					 (Q)--(N);
				\foreach \i/\g  in {A'/110,B'/90,C'/90,D'/120,A/-90,B/-90,C/-90,D/-90,P/90,Q/180,M/-90,N/-30}{\draw[fill=blue](\i) circle (1.5pt)
					 ($(\i)+(\g:5mm)$) node[scale=1]{$\i$};}
				\coordinate (Z) at  ($(D')+(0,2)$);
				 \coordinate (Y) at ($(C)+(2,0)$);
				\coordinate (X)  at ($(A)-(2,2)$);
				\draw[->]  (D')--(Z);
				\draw[->] (A)--(X); 
				 \draw[->] (C)--(Y);
				\draw (X)  node[scale=1,left]{$x$};
				\draw (Y)  node[scale=1,right]{$y$};
				\draw (Z) node[scale=1,above]{$z$}; 
			\end{tikzpicture}
	\end{center}
	 \shortans{-1}
	 \loigiai{
		\begin{center}
			\begin{tikzpicture}[scale=0.8,font=\footnotesize,line join=round,line cap=round,>=stealth]
					\coordinate (D) at (0,0); 
					\coordinate  (A) at (-2,-2);
					\coordinate (C) at  (6,0);
					\coordinate  (B) at ($(A)+(C)-(D)$);
					\foreach  \i in {A,B,C,D}{\coordinate (\i') at ($(\i)+(0,4)$);}
					\draw  (A')--(B')--(C')--(D')--cycle;
					 \draw (A)--(A')
						 (C)--(C')
						(B)--(B') 
						(A)--(B)--(C); 
					 \draw[dashed,thin](A)--(D)--(D')
						(C)--(D); 
					\coordinate (P) at  ($(D')!0.5!(C')$);
					 \coordinate (Q) at ($(D')!0.5!(D)$);
					 \coordinate (M) at ($(A)!0.5!(B)$);
					\coordinate (N) at ($(B)!0.5!(C)$); 
					\draw[dashed,thin](P)--(Q)--(M)--(N)--(P)--(M) 
						 (Q)--(N);
					\foreach \i/\g in {A'/110,B'/90,C'/90,D'/120,A/-90,B/-90,C/-90,D/-90,P/90,Q/180,M/-90,N/-30}{\draw[fill=blue](\i) circle (1.5pt) 
						($(\i)+(\g:5mm)$)  node[scale=1]{$\i$};}
					\coordinate (Z) at ($(D')+(0,2)$); 
					\coordinate  (Y) at ($(C)+(2,0)$);
					\coordinate (X) at  ($(A)-(2,2)$);
					 \draw[->] (D')--(Z);
					 \draw[->] (A)--(X);
					 \draw[->] (C)--(Y);
					\draw  (X) node[scale=1,left]{$x$};
					\draw (Y) node[scale=1,right]{$y$}; 
					\draw (Z) node[scale=1,above]{$z$}; 
				\end{tikzpicture}
		\end{center}
		Gắn hệ trục toạ độ như hình vẽ ta có $D\left(0;0;0\right),A\left(1;0;0\right)$, $B\left(1;1;0\right),C\left(0;1;0\right)$, $A'\left(1;0;1\right)$,  $B'\left(1;1;1\right)$, $C'\left(0;1;1\right),D'\left(0;0;1\right)$.\\
		Vì $M$, $N$, $P$, $Q$ lần  lượt là trung điểm của các cạnh $AB$, $BC$, $C'D'$ và $DD'$ nên $M\left(1;\dfrac{1}{2};0\right)$, $N\left(\dfrac{1}{2};1;0\right)$, $P\left(0;\dfrac{1}{2};1\right)$, $Q\left(0;0;\dfrac{1}{2} \right)$.\\
		Suy ra $\overrightarrow{MN}=\left(-\dfrac{1}{2};\dfrac{1}{2};0\right)$, $\overrightarrow{MP}=\left(-1;0;1\right)$, $\overrightarrow{MQ}=\left(-1;-\dfrac{1}{2};\dfrac{1}{2} \right)$. 
	} 
\end{bt}

% nguồn: Lop_12_Hinh_hoc_Chuong_2
\begin{bt}%[2H2H2-3]%[ma:f084c567]
	Cho hình chóp $S.ABC$ có đáy $ABC$ là tam giác vuông cân tại $A$ và $AB=a$, cạnh bên $SA$ vuông góc với mặt phẳng đáy và $SA=a$. Gọi $D$, $E$, $F$ lần lượt là điểm đối xứng của $A$ qua $C$, của $S$ qua $B$ và của $A$ qua mặt phẳng $\left(SBC\right)$. Khi $a=1$ thì tung độ của vectơ $\overrightarrow{AD}$ bằng bao nhiêu? 
	\begin{center}
		\begin{tikzpicture}[scale=1,font=\footnotesize,line join=round,line cap=round,>=stealth]
				\coordinate (A) at (0,0); 
				 \coordinate (S) at ($(A)+(0,4)$);
				\coordinate (D) at  (6,0);
				 \coordinate (C) at (2.7,0);
				\coordinate (B) at  (2,-1);
				\coordinate  (E) at (3,-3.6);
				\coordinate (F)  at (3,2);
				\coordinate (H)  at ($(A)!0.5!(F)$);
				\coordinate  (M) at ($(B)!0.5!(C)$);
				\coordinate (N)  at ($(B)!0.5!(S)$);
				\coordinate (I) at ($(A)!0.62!(F)$); 
				 \draw[dashed] (A)--(D)
					 (B)--(C)
					(A)--(H)--(C)--(M)--(H) 
					(I)--(C); 
				\draw (S)--(A)--(E)--(D)--(F)--(E)--(S) 
					(A)--(B) 
					 (F)--(H)--(S)--(B)
					(H)--(N) 
					 (S)--(I);
				\foreach \i/\g in  {C/60,A/180,B/-180,D/-30,E/-90,F/90,S/120,H/60}{\draw[fill=blue](\i) circle (1.5pt)
					($(\i)+(\g:5mm)$) node[scale=1]{$\i$};} 
				\pic[draw,thin,angle radius=2mm] {right angle =  A--H--M};
				\pic[draw,thin,angle radius=2mm] {right  angle = B--A--S};
				 \pic[draw,thin,angle radius=2mm] {right angle = C--A--S};
				\coordinate (Z)  at ($(S)+(0,1)$);
				\coordinate (Y)  at ($(A)!1.2!(D)$);
				\coordinate (X) at ($(A)!1.2!(B)$); 
				\draw[->]  (A)--(Z);
				 \draw[->] (B)--(X);
				 \draw[->] (D)--(Y);
				\draw  (X) node[scale=1,right]{$x$};
				\draw (Y) node[scale=1,above]{$y$}; 
				 \draw (Z) node[scale=1,left]{$z$};
			\end{tikzpicture}
	\end{center}
	 \shortans{2}
	 \loigiai{
		\begin{center}
			\begin{tikzpicture}[scale=1,font=\footnotesize,line join=round,line cap=round,>=stealth]
					 \coordinate (A) at (0,0);
					\coordinate (S) at ($(A)+(0,4)$); 
					\coordinate (D)  at (6,0);
					\coordinate (C) at  (2.7,0);
					\coordinate (B) at (2,-1); 
					\coordinate (E) at (3,-3.6); 
					\coordinate (F) at (3,2); 
					 \coordinate (H) at ($(A)!0.5!(F)$);
					\coordinate  (M) at ($(B)!0.5!(C)$);
					 \coordinate (N) at ($(B)!0.5!(S)$);
					\coordinate  (I) at ($(A)!0.62!(F)$);
					 \draw[dashed] (A)--(D)
						 (B)--(C)
						(A)--(H)--(C)--(M)--(H) 
						(I)--(C); 
					 \draw (S)--(A)--(E)--(D)--(F)--(E)--(S)
						 (A)--(B)
						 (F)--(H)--(S)--(B)
						 (H)--(N)
						(S)--(I); 
					\foreach \i/\g in  {C/60,A/180,B/-180,D/-30,E/-90,F/90,S/120,H/60}{\draw[fill=blue](\i) circle (1.5pt)
						($(\i)+(\g:5mm)$) node[scale=1]{$\i$};} 
					\pic[draw,thin,angle  radius=2mm] {right angle = A--H--M};
					\pic[draw,thin,angle radius=2mm] {right angle = B--A--S}; 
					\pic[draw,thin,angle  radius=2mm] {right angle = C--A--S};
					\coordinate (Z) at ($(S)+(0,1)$); 
					\coordinate  (Y) at ($(A)!1.2!(D)$);
					 \coordinate (X) at ($(A)!1.2!(B)$);
					\draw[->] (A)--(Z); 
					\draw[->]  (B)--(X);
					\draw[->] (D)--(Y); 
					\draw (X)  node[scale=1,right]{$x$};
					 \draw (Y) node[scale=1,above]{$y$};
					\draw  (Z) node[scale=1,left]{$z$};
				\end{tikzpicture}
		\end{center}
		Chọn hệ trục tọa độ $Axyz$ như hình vẽ với độ dài vectơ đơn vị  bằng $a$, ta được tọa độ các điểm như sau $A\left(0;0;0\right)$, $B\left(1;0;0\right)$, $C\left(0;1;0\right)$, $S\left(0; 0; 1\right)$ $\Rightarrow D\left(0;2;0\right)$, $E\left(2;0;-1\right)$.\\
		Gọi $H$  là trọng tâm tam giác $SBC$ $\Rightarrow H\left(\dfrac{1}{3};\dfrac{1}{3};\dfrac{1}{3} \right)$. Vì tam giác $SBC$ là tam giác đều nên $H$ cũng là trực tâm của tam giác $SBC$ $\Rightarrow AH\perp \left(SBC\right)$ $\Rightarrow F$ là điểm đối xứng với $A$ qua $H$.\\
		$\Rightarrow F\left(\dfrac{2}{3};\dfrac{2}{3};\dfrac{2}{3} \right)$. Ta tính  được $\overrightarrow{AD}=\left(0;2;0\right)$, $\overrightarrow{AE}=\left(2;0;-1\right)$, $\overrightarrow{AF}=\left(\dfrac{2}{3};\dfrac{2}{3};\dfrac{2}{3} \right)$.
	 }
\end{bt}

% nguồn: Lop_12_Hinh_hoc_Chuong_2
\begin{bt}%[2H2H2-3]%[ma:8d7c16c6]
	Trong không gian $Oxyz$, cho hai điểm $A(1;2;3)$, $B(-2;-4;9)$. Điểm $M$ thuộc đoạn $AB$ sao cho $MA=2MB$. Bình phương độ dài đoạn thẳng $OM$ bằng  bao nhiêu?
	\shortans{54} 
	\loigiai{ 
		\begin{center}
			\begin{tikzpicture}[scale=1,font=\footnotesize,line join=round,
						line  cap=round,>=stealth]
					\draw (0:0) coordinate (A)--++(0:4) coordinate (M)--++(0:2) coordinate (B); 
					 \foreach \x/\g in {A/90,B/90,M/90}
					\fill (\x) circle (1pt) 
						($(\x)+(\g:3mm)$) node{$\x$}; 
				\end{tikzpicture}
		\end{center}
		Gọi $M(x;y;z)$. Vì điểm $M$ thuộc đoạn $AB$ sao cho $MA=2MB$ nên 
		 $$
		\overrightarrow{AB}=3  \overrightarrow{MB} \Leftrightarrow\heva{&3 \cdot (-2-x)=-3\\&3 \cdot (-4-y)=-6\\&3\cdot(9-z)=6} \Leftrightarrow\heva{&x=-1 \\&y=-2 \\& z=7.}$$
		Suy ra $M(-1;-2;7)$.\\  Từ đó có $OM=\sqrt{1+4+49}=\sqrt{54} \Rightarrow OM^2=54$.
	} 
\end{bt}

% nguồn: Lop_12_Hinh_hoc_Chuong_2
\begin{bt}%[2H2H2-3]%[ma:3086d314]
	Cho lăng trụ đứng $ABCD.A'B'C'D'$ có đáy $ABCD$ là hình thoi tâm $O$ cạnh $a$, $AA'=a$ và $\widehat{ABC}=60^{\circ}$. Gọi $M$ là trung điểm của cạnh $AA'$. Chọn hệ trục tọa độ như hình vẽ.  Vectơ $\overrightarrow{OC'}$ có toạ độ là $\left(m; n; p\right)$. Khi $a=1$ hãy tính giá trị biểu thức $T=2m+n+p$.
	\begin{center}
		\begin{tikzpicture}[scale=0.8,font=\footnotesize,line join=round,line cap=round,>=stealth]
				\coordinate  (A) at (0,0);
				\coordinate (B) at  (-2,-2);
				\coordinate (D)  at (6,0);
				 \coordinate (C) at ($(B)+(D)-(A)$);
				\foreach \i in {A,B,C,D}{\coordinate (\i')  at ($(\i)+(0,4)$);}
				\coordinate  (M) at ($(A)!0.5!(A')$);
				\coordinate (O) at ($(A)!0.5!(C)$); 
				 \draw (A')--(B')--(C')--(D')--cycle;
				\draw  (B)--(B')
					 (C)--(C')--(D)--(D')
					 (B)--(C)--(D);
				\draw[dashed,thin](B)--(A)--(A') 
					(A)--(D) 
					(M)--(O)--(A) 
					(O)--(B); 
				\foreach  \i/\g in {A'/110,B'/90,C'/90,D'/90,A/-90,B/-90,C/-90,D/-90,M/180,O/-90}{\draw[fill=blue](\i) circle (1.5pt)
					($(\i)+(\g:5mm)$)  node[scale=1]{$\i$};}
				 \coordinate (Z) at ($(O)+(0,6)$);
				\coordinate (Y)  at ($(O)!1.5!(D)$);
				\coordinate (X) at  ($(O)!1.8!(C)$);
				\draw[->]  (O)--(Z);
				\draw[->]  (O)--(X);
				\draw[->]  (O)--(Y);
				\draw (X) node[scale=1,left]{$x$}; 
				\draw  (Y) node[scale=1,above]{$y$};
				\draw (Z) node[scale=1,left]{$z$}; 
			\end{tikzpicture}
	\end{center}
	 \shortans{2}
	\loigiai{ 
		Dựng  hệ trục tọa độ như hình vẽ.
		\begin{center}
			\begin{tikzpicture}[scale=0.8,font=\footnotesize,line join=round,line cap=round,>=stealth]
					\coordinate  (A) at (0,0);
					\coordinate  (B) at (-2,-2);
					 \coordinate (D) at (6,0);
					\coordinate (C) at ($(B)+(D)-(A)$); 
					\foreach \i in {A,B,C,D}{\coordinate (\i') at  ($(\i)+(0,4)$);}
					\coordinate (M) at  ($(A)!0.5!(A')$);
					\coordinate (O) at  ($(A)!0.5!(C)$);
					 \draw (A')--(B')--(C')--(D')--cycle;
					\draw (B)--(B') 
						 (C)--(C')--(D)--(D')
						 (B)--(C)--(D);
					 \draw[dashed,thin](B)--(A)--(A')
						(A)--(D) 
						 (M)--(O)--(A)
						 (O)--(B);
					\foreach \i/\g  in {A'/110,B'/90,C'/90,D'/90,A/-90,B/-90,C/-90,D/-90,M/180,O/-90}{\draw[fill=blue](\i) circle (1.5pt)
						($(\i)+(\g:5mm)$) node[scale=1]{$\i$};} 
					\coordinate (Z) at ($(O)+(0,6)$); 
					\coordinate  (Y) at ($(O)!1.5!(D)$);
					\coordinate (X) at  ($(O)!1.8!(C)$);
					 \draw[->] (O)--(Z);
					\draw[->]  (O)--(X);
					\draw[->]  (O)--(Y);
				\end{tikzpicture}
		\end{center}
		Gọi  $O$ là giao điểm của $AC$ và $BD$. Chọn $a=1$.
		Dễ dàng  tính được $O\left(0;0;0\right)$, $M\left(-\dfrac{1}{2};0;\dfrac{1}{2} \right)$, $D\left(0;\dfrac{\sqrt{3}}{2};0\right)$ và $C'\left(\dfrac{1}{2};0;1\right)$.\\
		Suy ra $\overrightarrow{OM}=\left(-\dfrac{1}{2};0;\dfrac{1}{2} \right)$, $\overrightarrow{C'D}=\left(-\dfrac{1}{2};\dfrac{\sqrt{3}}{2};-1\right)$ và $\overrightarrow{OC'}=\left(\dfrac{1}{2};0;1\right)$.\\ 
		Khi đó $\heva{m&=\dfrac{1}{2} \\ n&=0 \\ p&=1} \Rightarrow T=2m+n+p=2\cdot  \dfrac{1}{2}+0+1=2$.
	} 
\end{bt}

% nguồn: Lop_12_Hinh_hoc_Chuong_2
\begin{bt}%[2H2H2-3]%[ma:983792b3]
	Cho lăng trụ $ABC.A'B'C'$ có tam giác $ABC$ đều cạnh bằng $2a$. Hình chiếu vuông góc của $A'$ lên mặt phẳng $\left(ABC\right)$ trùng với trung điểm $M$ của $BC$ và góc giữa $A'B$ và mặt phẳng $(ABC)$ là $60^{\circ}$. Khi $a=1$ thì tung  độ của $\overrightarrow{AB'}$ bằng bao nhiêu?
	\begin{center}
		\begin{tikzpicture}[scale=1,font=\footnotesize,line join=round,line cap=round,>=stealth]
				\coordinate (M)  at (0,0);
				\coordinate (A') at  ($(M)+(0,4)$);
				\coordinate (C) at (2,1); 
				\coordinate (A)  at (-2.5,0.5);
				\coordinate  (B) at ($(C)!2!(M)$);
				 \coordinate (C') at ($(A')+(C)-(A)$);
				\coordinate (B') at  ($(A')+(B)-(A)$);
				\filldraw[green](A)--(C)--(B')--cycle; 
				\draw[dashed]  (A)--(M)--(A')
					 (A)--(C);
				\draw (A)--(B)--(C)--(C')--(A')--(B')--(B)--(A')--(A) 
					(B')--(C'); 
				\foreach \i/\g in {A/-160,B/-90,C/-20,A'/150,B'/-30,C'/-10}{\draw[fill=blue](\i) circle (1.5pt) 
					($(\i)+(\g:5mm)$) node[scale=1]{$\i$};} 
				\pic[draw,thin,angle radius=2mm] {right angle =  C--M--A'};
				\coordinate (Z)  at ($(A')+(0,1)$);
				\coordinate  (Y) at ($(M)!1.2!(B)$);
				\coordinate (X)  at ($(M)!1.2!(A)$);
				\draw[->] (A')--(Z); 
				 \draw[->] (B)--(Y);
				\draw[->] (A)--(X); 
				\draw (X) node[scale=1,right]{$x$}; 
				\draw (Y)  node[scale=1,above]{$y$};
				 \draw (Z) node[scale=1,left]{$z$};
				\draw (M) node[scale=1,below]{$M\equiv O$}; 
			\end{tikzpicture}
	\end{center}
	\shortans{1} 
	\loigiai{ 
		\begin{center}
			\begin{tikzpicture}[scale=1,font=\footnotesize,line join=round,line cap=round,>=stealth]
					\coordinate (M)  at (0,0);
					\coordinate (A') at ($(M)+(0,4)$); 
					\coordinate (C) at  (2,1);
					 \coordinate (A) at (-2.5,0.5);
					\coordinate  (B) at ($(C)!2!(M)$);
					\coordinate (C')  at ($(A')+(C)-(A)$);
					\coordinate (B') at  ($(A')+(B)-(A)$);
					\filldraw[green](A)--(C)--(B')--cycle; 
					 \draw[dashed] (A)--(M)--(A')
						 (A)--(C);
					 \draw (A)--(B)--(C)--(C')--(A')--(B')--(B)--(A')--(A)
						(B')--(C'); 
					\foreach \i/\g  in {A/-160,B/-90,C/-20,A'/150,B'/-30,C'/-10}{\draw[fill=blue](\i) circle (1.5pt)
						($(\i)+(\g:5mm)$) node[scale=1]{$\i$};} 
					\pic[draw,thin,angle radius=2mm]  {right angle = C--M--A'};
					 \coordinate (Z) at ($(A')+(0,1)$);
					\coordinate (Y) at  ($(M)!1.2!(B)$);
					\coordinate (X) at  ($(M)!1.2!(A)$);
					 \draw[->] (A')--(Z);
					\draw[->]  (B)--(Y);
					\draw[->] (A)--(X); 
					 \draw (X) node[scale=1,right]{$x$};
					 \draw (Y) node[scale=1,above]{$y$};
					 \draw (Z) node[scale=1,left]{$z$};
					\draw  (M) node[scale=1,below]{$M\equiv O$};
				\end{tikzpicture}
		\end{center}
		Góc giữa $A'B$ và mặt phẳng $\left(ABC\right)$  là góc $\widehat{A'BM}=60^{\circ}$.\\
		Ta  có $A'M=BM\cdot \tan 60^{\circ}=a\sqrt{3}$; $AM=AB\cdot \sin 60^{\circ}=2a\cdot \dfrac{\sqrt{3}}{2}=a\sqrt{3}$.\\
		Đặt hệ trục toạ $Oxyz$ như hình  vẽ. Khi đó $A\left(\sqrt{3};0;0\right)$, $B\left(0;1;0\right)$, $C\left(0;-1;0\right)$,
		$A'\left(0;0;\sqrt{3} \right)$.  \\
		Gọi $B'\left(x_0;y_0;z_0  \right)$, $\overrightarrow{A'B'}=\left(x_0;y_0;z_0-\sqrt{3} \right)$, $\overrightarrow{AB}=\left(-\sqrt{3};1;0\right)$.\\
		Vì $\overrightarrow{A'B'}=\overrightarrow{AB}$ nên  $\heva{x_0&=-\sqrt{3} \\ y_0&=1 \\ z_0&=\sqrt{3}} $ $\Rightarrow B'\left(-\sqrt{3};1;\sqrt{3} \right)\Rightarrow \overrightarrow{AB'}=\left(-2\sqrt{3};1;\sqrt{3} \right)$.
	} 
\end{bt}

% nguồn: Lop_12_Hinh_hoc_Chuong_2
\begin{bt}%[2H2H2-3]%[ma:4b8439f0]
	Trong không gian $Oxyz$, cho tứ diện $ABCD$ với $A(1;-4;2)$, $B(2;1;-3)$, $C(3;0;-2)$ và $D(2;-5;-1)$. Hoành độ điểm $G$ thỏa mãn $\overrightarrow{GA}+\overrightarrow{GB}+\overrightarrow{GC}+\overrightarrow{GD}=\overrightarrow{0}$ là  bao nhiêu?
	\shortans{2} 
	\loigiai{ 
		Tọa độ điểm $G$ thỏa mãn $\heva{&x_G=\dfrac{x_A+x_B+x_C+x_D}{4}=\dfrac{1+2+3+2}{4}=2 \\& y_G=\dfrac{y_A+y_B+y_C+y_D}{4}=\dfrac{-4+1+0+(-5)}{4}=-2  \\& z_G=\dfrac{z_A+z_B+z_C+z_D}{4}=\dfrac{2+(-3)+(-2)+(-1)}{4}=-1} \Rightarrow G(2;-2;-1)$.
	} 
\end{bt}

% nguồn: Lop_12_Hinh_hoc_Chuong_2
\begin{bt}%[2H2V2-3]%[ma:074ba77b]
	Trong không gian $Oxyz$, cho $3$ điểm $A(5;-2;0)$, $B(4;5;-2)$ và $C(0;3;2)$. Điểm $M$ di chuyển trên trục $Ox$. Đặt $Q=2\left| \overrightarrow{MA}+\overrightarrow{MB}+\overrightarrow{MC}\right|+3\left|\overrightarrow{MB}+\overrightarrow{MC}\right|$. Biết giá trị nhỏ nhất của $Q$ có dạng $a\sqrt{b}$ trong đó $a, b  \in \mathbb{N}$ và $b$ là số nguyên tố. Tính $a+b$.
	 \shortans{43}
	 \loigiai{
		Ta  có
		$$Q=2\left| \overrightarrow{MA}+\overrightarrow{MB}+\overrightarrow{MC}\right|+3\left|\overrightarrow{MB}+\overrightarrow{MC}\right|=2\left|3\overrightarrow{MG}+\overrightarrow{GA}+\overrightarrow{GB}+\overrightarrow{GC}\right|+3\left|2\overrightarrow{MI}+\overrightarrow{IB}+\overrightarrow{IC}\right|.$$ 
		Với $G(3;2;0)$ là trọng tâm của tam giác $ABC$ và $I(2;4;0)$ là  trung điểm $BC$, ta có
		$$Q=2\left|3  \overrightarrow{MG}\right| +3\left|2 \overrightarrow{MI}\right|=6(MG+MI).$$
		Do $G$ và $I$ nằm cùng phía so với $Ox$ nên ta có $G'(3;-2;0)$ là điểm đối xứng của $G$  qua $Ox$.\\
		Khi  đó $Q=2\left|3\overrightarrow{MG}\right|+3\left|2\overrightarrow{MI}\right|=6(MG+MI)=6\left(MG'+MI\right) \geq 6G'I=6\sqrt{37}$.\\
		Đẳng thức  xảy ra khi $M$ là giao điểm của $G'I$ và $Ox$.
	 }
\end{bt}

% nguồn: Lop_12_Hinh_hoc_Chuong_2
\begin{bt}%[2H2V2-3]%[ma:dad17b9c]
	Trong không gian $Oxyz$, cho ba điểm $A(-2;3;1)$, $B(2;1;0)$, $C(-3;-1;1)$. Gọi $D(a;b;c)$ là điểm sao cho $ABCD$ là hình thang có  cạnh đáy $AD$ và diện tích hình thang $ABCD$ bằng $4$ lần diện tích tam giác $ABC$. Tính $a+b+c$.
	\shortans{-16} 
	 \loigiai{
		\begin{center}
			\begin{tikzpicture}[scale=1,font=\footnotesize,line join=round,
						line cap=round,>=stealth] 
					\draw (0:0) coordinate (D)--++(180:4.5) coordinate (A)--++(70:2.6) coordinate  (B)--++(0:1.5) coordinate (C)--cycle;
					\foreach  \x/\g in {A/180,B/180,C/0,D/0}
					\fill  (\x) circle (1pt)
						($(\x)+(\g:3mm)$)  node{$\x$};
				\end{tikzpicture}
		\end{center}
		Ta có 
		 \allowdisplaybreaks
		\begin{eqnarray*}
			S_{ABCD}=4S_{ABC} &\Leftrightarrow& \dfrac{1}{2}\cdot \mathrm{d}(BC,AD)\cdot (BC+AD)=4 \cdot \dfrac{1}{2}\cdot \mathrm{d}(BC,AD)\cdot BC\\ 
			 &\Leftrightarrow& BC+AD=4BC\\
			 &\Leftrightarrow& AD=3BC.
		\end{eqnarray*}
		Do $ABCD$ là hình thang có đáy $AD$  nên
		$$ \overrightarrow{AD}=3\overrightarrow{BC}\Leftrightarrow\heva{&a+2=-15\\&b-3=-6\\&c-1=3} \Leftrightarrow\heva{&a=-17  \\& b=-3 \\&c=4} \Rightarrow a+b+c=-16.$$
	} 
\end{bt}

% nguồn: Lop_12_Hinh_hoc_Chuong_2
\begin{bt}%[2H2V2-3]%[ma:0196497c]
	Trong không gian $Oxyz$, cho hai điểm $A(2;-1;-5)$, $B(-4;2;1)$. Xét $M$ là điểm thay đổi thỏa mãn điều kiện $\left|\overrightarrow{MA}+2\overrightarrow{MB}\right|=9$. Độ dài đoạn thẳng  $OM$ lớn nhất bằng bao nhiêu (kết quả làm tròn đến hàng phần trăm)?
	\shortans{5{,45}}%CẦN KIỂM TRA ĐÁP ÁN 
	 \loigiai{
		Trước  hết, ta tìm tọa độ điểm $I$ sao cho $\overrightarrow{IA}+2\overrightarrow{IB}=\overrightarrow{0}$.
		$$\overrightarrow{IA}+2\overrightarrow{IB}=\overrightarrow{0}\Leftrightarrow \overrightarrow{OA}+2\overrightarrow{OB}=3\overrightarrow{OI}\Leftrightarrow  \overrightarrow{OI} =\dfrac{1}{3}\overrightarrow{OA}+\dfrac{2}{3}\overrightarrow{OB}.$$
		Suy ra $\heva{&x_I=\dfrac{x_A+2x_B}{3}=-2\\&y_I=\dfrac{y_A+2y_B}{3}=1\\& z_I=\dfrac{z_A+2z_B}{3}=-1. }$ Do đó,  $I(-2;1;-1)$.\\
		Hơn nữa  $\overrightarrow{MA}+2\overrightarrow{MB}=3\overrightarrow{MI}$ nên $MI=\left|\vv{MI}\right|=\dfrac{1}{3}\left|\overrightarrow{MA}+2\overrightarrow{MB}\right|=3$.\\
		Vậy tập hợp các  điểm $M$ là mặt cầu tâm $I(-2;1;-1)$, bán kính $R=3$.\\
		Suy  ra độ dài lớn nhất của đoạn $OM=OI+R=\sqrt{6}+3\approx 5{,45}$.
	 }
\end{bt}


%%%%%%%%%%%%%%%% TỰ LUẬN %%%%%%%%%%%%%%%%

% nguồn: Lop_12_Hinh_hoc_Chuong_2
\begin{bt}%[2H2H2-3]%[ma:ec70a4bc]
	Trong không gian $Oxyz$, cho hình hộp chữ nhật $ABCD.A'B'C'D'$ có đỉnh $A$ trùng với gốc $O$, các vectơ  $\overrightarrow{AB}$, $\overrightarrow{AD}$, $\overrightarrow{AA'}$ theo thứ tự cùng hướng với $\overrightarrow{i}$, $\overrightarrow{j}$, $\overrightarrow{k}$ và có $AB=8$, $AD=6$, $AA'=4$. Tìm tọa độ các vectơ $\overrightarrow{AB}$, $\overrightarrow{AC}$, $\overrightarrow{AC'}$ và $\overrightarrow{AM}$ với $M$ là trung điểm của cạnh $C'D'$.
	\loigiai{ 
		\begin{center}
			\begin{tikzpicture}[scale=0.55,font=\footnotesize,line join=round,
						line  cap=round,>=stealth]
					\def\B{6} 
					\def\C{8} 
					\def\h{4} 
					\path (0:0) coordinate  (A)
						++(0:\B) coordinate  (D)
						 ++(-130:\C) coordinate (C)
						($(A)+(C)-(D)$) coordinate (B) 
						($(A)+(90:\h)$) coordinate  (A')
						 ($(D)+(90:\h)$) coordinate (D')
						 ($(B)+(90:\h)$) coordinate (B')
						($(C)+(90:\h)$) coordinate  (C')
						(A')++(90:1)  coordinate (z)
						(B)++(-130:0.7)  coordinate (x)
						(D)++(0:1) coordinate  (y)
						($(C')!.5!(D')$)  coordinate (M);
					\draw[dashed]  (B)--(A)--(D)
						(A)--(A'); 
					 \draw (D)--(D')
						 (C)--(C')
						(B)--(B') 
						(B)--(C)--(D) 
						(B')--(C')--(D')--(A')--cycle; 
					 \draw[->] (A')--(z);
					\draw[->]  (B)--(x);
					\draw[->]  (D)--(y);
					\foreach \x/\g in  {B/180,C/0,D/-90,A/180,B'/180,C'/0,D'/0,A'/180,x/180,y/90,z/0,M/135}
					\fill (\x) circle  (1pt)
						($(\g:4mm)+(\x)$)  node {$\x$};
				\end{tikzpicture}
		\end{center}
		Để tìm tọa độ của vectơ $\overrightarrow{AB}$ ta cần biểu diễn $\overrightarrow{AB}$ theo ba vectơ $\overrightarrow{i}$,  $\overrightarrow{j}$, $\overrightarrow{k}$.\\
		Do $\overrightarrow{AB}$ cùng hướng  với $\overrightarrow{i}$ và $|\overrightarrow{AB}|=AB=8=8|\overrightarrow{i}|$ nên $\overrightarrow{AB}=8\overrightarrow{i}$ hay $\overrightarrow{AB}=8\overrightarrow{i}+0\overrightarrow{j}+0\overrightarrow{k}$.\\
		Tương tự,  ta cũng có $\overrightarrow{AD}=0\overrightarrow{i}+6\overrightarrow{j}+0\overrightarrow{k}$, $\overrightarrow{AA'}=0\overrightarrow{i}+0\overrightarrow{j}+4\overrightarrow{k}$.\\
		Trong hình bình hành $ABCD$ ta  có $\overrightarrow{AC}=\overrightarrow{AB}+\overrightarrow{AD}=8\overrightarrow{i}+6\overrightarrow{j}+0\overrightarrow{k}$.\\
		Trong hình bình hành $AA'C'C$ ta có  $\overrightarrow{AC'}=\overrightarrow{AC}+\overrightarrow{AA'}=8\overrightarrow{i}+6\overrightarrow{j}+4\overrightarrow{k}$.\\
		Suy ra  $\overrightarrow{AB}=(8;0;0)$, $\overrightarrow{AC}=(8;6;0)$, $\overrightarrow{AC'}=(8;6;4)$.\\
		Vì  $\overrightarrow{AM}=\dfrac{1}{2}\left(\overrightarrow{AC'}+\overrightarrow{AD'}\right)=\dfrac{1}{2}\left(\overrightarrow{AC'}+\overrightarrow{AD}+\overrightarrow{AA'}\right)=\dfrac{1}{2}\left(8\overrightarrow{i}+6\overrightarrow{j}+4\overrightarrow{k}+6\overrightarrow{j}+4\overrightarrow{k}\right)=4\overrightarrow{i}+6\overrightarrow{j}+4\overrightarrow{k}$.\\
		Suy ra  $\overrightarrow{AM}=(4;6;4)$.
	} 
\end{bt}

% nguồn: Lop_12_Hinh_hoc_Chuong_2
\begin{bt}%[2H2H2-3]%[ma:03a153f2]
	 Trong không gian $Oxyz$, cho ba vectơ $\overrightarrow{p}=(3;-2;1)$, $\overrightarrow{q}=(6;-4;2)$, $\overrightarrow{r}=(2;1;-3)$.
	\begin{enumerate}
		\item  Tìm toạ độ của vectơ $\overrightarrow{c}=2\overrightarrow{p}-3\overrightarrow{q}+\overrightarrow{r}$.
		\item Tìm hai vectơ cùng phương  trong các vectơ đã cho.
	\end{enumerate}
	 \loigiai{
		\begin{enumerate}
			\item  Ta có $2\overrightarrow{p} = (6;-4;2)$, $-3\overrightarrow{q}=(-18;12;-6)$, $\overrightarrow{r}=(2;1;-3)$. \\
			Suy ra  $\overrightarrow{c}=(6-18+2;-4+12+1;2-6-3)=(-10;9;-7)$.
			\item Ta thấy  $\overrightarrow{q}=(6;-4;2)=2(3;-2;1)=2\overrightarrow{p}$.\\
			Vậy  hai vectơ $\overrightarrow{p}, \overrightarrow{q}$ cùng phương.
		\end{enumerate}
	 }
\end{bt}

% nguồn: Lop_12_Hinh_hoc_Chuong_2
\begin{bt}%[2H2H2-3]%[ma:3acfe4d2]
	Cho $\overrightarrow{a}=(2;-1;5)$, $\overrightarrow{b}=(0;3;-3)$, $\overrightarrow{c}=(1;4;-2)$.  Tìm toạ độ của vectơ $\overrightarrow{d} =2\overrightarrow{a}-\dfrac{1}{5}\overrightarrow{b}+3\overrightarrow{c}$.
	\loigiai{ 
		Ta có $2\overrightarrow{a} = (4;-2;10)$, $-\dfrac{1}{5}\overrightarrow{b}  = \left(0;-\dfrac{3}{5};\dfrac{3}{5}\right)$, $3\overrightarrow{c} = (3;12;-6)$.\\
	Nên  $\overrightarrow{d} = \left(4+0+3;-2-\dfrac{3}{5}+12;10+\dfrac{3}{5}-6\right)=\left(7;\dfrac{47}{5};\dfrac{23}{5}\right)$.}
\end{bt}

% nguồn: Lop_12_Hinh_hoc_Chuong_2
\begin{bt}%[2H2H2-3]%[ma:089aacab]
	Trong không gian $Oxyz$, cho ba  điểm không thẳng hàng $A(2;-1;4)$, $B(3;5;-1)$, $C(-1;1;2)$.
	\begin{enumerate}
		 \item Tìm toạ độ của $\overrightarrow{AB}$.
		 \item Tìm tọa độ điểm $D$ sao cho $ABCD$ là hình bình hành.
	\end{enumerate}
	\loigiai{ 
		\begin{enumerate}
			 \item Vì $A(2;-1;4)$ và $B(3;5;-1)$ nên $\overrightarrow{AB}=(3-2;5-(-1);-1-4)$ nên $\overrightarrow{AB}=(1;6;-5)$.
			\item Do ba điểm $A$, $B$, $C$ không  thẳng hàng nên để $ABCD$ là hình bình hành thì điểm $D$ phải thoả mãn điều kiện $\overrightarrow{AB}=\overrightarrow{DC}$.\\
			Gọi $\left(x_D;y_D;z_D\right)$ là toạ độ điểm $D$.\\ 
			Ta có  $\overrightarrow{DC}=\left(-1-x_D;1-y_D;2-z_D\right)$.\\
			Suy ra  $\overrightarrow{AB}=\overrightarrow{DC}\Leftrightarrow\heva{&-1-x_D=1\\&1-y_D=6\\&2-z_D=-5}\Leftrightarrow\heva{&x_D=-2\\&y_D=-5\\&z_D=7.}$\\
			Vậy $D(-2;-5;7)$ thì tứ giác $ABCD$ là hình  bình hành.
		\end{enumerate}
	} 
\end{bt}

% nguồn: Lop_12_Hinh_hoc_Chuong_2
\begin{bt}%[2H2H2-3]%[ma:dfd01429]
	Trong không gian  $Oxyz$, cho vectơ $\overrightarrow{u}=-2\overrightarrow{i}+3\overrightarrow{j}+\dfrac{3}{4}\overrightarrow{k}$ và vectơ $\overrightarrow{v} = \left(3;-\dfrac{5}{4};2\right)$.
	\begin{enumerate}
		\item Tìm toạ độ của $\overrightarrow{u}$. 
		\item Biểu diễn $\overrightarrow{v}$ theo các vectơ  đơn vị $\overrightarrow{i}$, $\overrightarrow{j}$, $\overrightarrow{k}$.
		\item Tìm toạ độ  của $\overrightarrow{a}=2\overrightarrow{u}+\dfrac{1}{3}\overrightarrow{v}$.
	\end{enumerate}
	 \loigiai{
		\begin{enumerate}
			\item Từ $\overrightarrow{u} = -2\overrightarrow{i} + 3\overrightarrow{j}  + \dfrac{3}{4}\overrightarrow{k}$. Suy ra $\overrightarrow{u} = \left(-2;3;\dfrac{3}{4}\right)$.
			\item Từ $\overrightarrow{v} =  \left(3;-\dfrac{5}{4};2\right)$. Suy ra $\overrightarrow{v}=3\overrightarrow{i}-\dfrac{5}{4}\overrightarrow{j}+2\overrightarrow{k}$.
			\item Ta có $2\overrightarrow{u}=\left(-4;6;\dfrac{3}{2}\right)$ và $\dfrac{1}{3}\overrightarrow{v}  = \left(1;-\dfrac{5}{12}; \dfrac{2}{3}\right)$.\\
			 Do đó $\overrightarrow{a} = 2\overrightarrow{u} + \dfrac{1}{3}\overrightarrow{v}=\left(-4+1;6-\dfrac{5}{12};\dfrac{3}{2}+\dfrac{2}{3}\right)=\left(-3;\dfrac{67}{12};\dfrac{13}{6}\right)$.
		\end{enumerate}
	 }
\end{bt}
